\chapter{Phrasen und interner Bau von Satzkonstituenten}\label{chap:Phrasen}

In dem vorangehenden Kapitel~\ref{chap:Konstituenz} haben wir gesehen, dass die unmittelbaren Satzkonstituenten komplex sein können. Um den inneren Aufbau jener Wortgruppen geht es in diesem Kapitel und zwar unabhängig davon, ob es sich um Satzkonstituenten handelt oder nicht. Der innere Bau bzw. die Ausbaufähigkeit einer Wortgruppe hängt von dem Wort ab, das das Zentrum der Wortgruppe bildet. Dieses Wort nennen wir\is{Kopf} \textit{Kopf} und die vom Kopf aus gebildete Einheit\is{Phrase} \textit{Phrase}. 

\begin{Definition}{Phrase, Kopf und Beziehungen}
    
\noindent Eine Phrase ist eine syntaktische Einheit, die von einem Kopf bestimmt wird. Eine Phrase kann aus mehreren Wörtern oder nur aus dem Kopfwort bestehen. Die Kategorie der Phrase entspricht der Kategorie ihres Kopfes. Der Kopf legt aufgrund seiner lexikalischen und kategorialen Eigenschaften bestimmte Formmerkmale der abhängigen Wörter fest (Rektion) und stimmt in veränderlichen Merkmalen mit ihnen überein (Kongruenz). Der Kopf legt die Stellung der von ihm abhängigen Wörter fest (Positionsbezug).     
\end{Definition}


Nach der Wortartzugehörigkeit des Kopfes unterscheidet man zwischen verschiedenen Phrasentypen. Diese werden wir in den sieben folgenden Abschnitten beschreiben: die Verbalphrase (VP) in Abschnitt~\ref{section:VP}, die Nominalphrase (NP) in Abschnitt~\ref{section:NP}, die Adjektivphrase (AP) in Abschnitt~\ref{section:AP}, die Präpositionalphrase (PP) in Abschnitt~\ref{section:PP}, die Subjunktionalphrase in Abschnitt~\ref{section:SubjunktP}, die Adverbphrase (AdvP) in Abschnitt~\ref{section:Adverbphrase} und abschließend die Adjunktionalphrase in Abschnitt~\ref{section:Adjunktionalphrase}. Stichpunktartig fassen wir dieses Kapitel im Abschnitt~\ref{section:Zfsg_Phrasen} zusammen. Der letzte Abschnitt~\ref{section:Aufgaben_Phrasen} bietet vier verschiedene Anwendungsaufgaben samt Lösungen zum Thema Phrasen. Zuvor gehen wir aber in dem folgenden Exkurs~\ref{ExkursPhrase_Schule} auf die schulische Relevanz des Phrasenkonzepts ein.  



\begin{Exkurs}{Phrasenbegriff und Schule}
\label{ExkursPhrase_Schule}
 
Insbesondere in der\is{Schulgrammatik} Schulgrammatik wurden und werden Phrasen bzw. Wortgruppen kaum behandelt. Wir benutzen wie die \citet[36, §11]{Duden2022} den etablierten Phrasenbegriff gleichbedeutend zu Wortgruppe. Da der Gruppenbegriff suggeriert, dass es sich um mehrere Wörter handelt, ziehen wir den abstrakteren Phrasenbegriff vor. Denn unter ihn kann ein einzelnes Wort wie \textit{er} genauso wie die Wortgruppe \textit{der kleine Junge, der heute Geburtstag hat} fallen.\footnote{Andere theoretische Grundannahmen wie die der x-bar-Theorie teilen wir nicht. Statt Nominalphrase könnten wir auch von einer Nominalkonstruktion sprechen. Wir verzichten auf die von \citet[136--137]{Himmelmann1997} vorgeschlagene Differenzierung zwischen mehr oder weniger selbständigen, nebeneinander stehenden (nur semantisch oder morphologisch zusammengehörenden) Wörtern als Wortgruppe und via Grammatikalisierung hierarchisch organisierten Phrasen. Eine solche Differenzierung ist vor allem dann nötig, wenn verschiedene Sprachen bzw. historische Sprachstufen miteinander verglichen werden.} 

Das schulische Manko liegt darin begründet, dass in der Schule auf der einen Seite Wortarten unterrichtet werden und auf der anderen Seite die Funktionen der Satzkonstituenten gegenüber dem Prädikat. Satzgliedfunktionen wie Subjekt und Objekt sowie der Begriff des Attributs für Teile von Satzkonstituenten geben aber keinen Aufschluss darüber, wie die entsprechenden Phrasen gebaut sind. Ohne die verschiedenen Möglichkeiten des Phrasenbaus zu thematisieren, können die formalen Unterschiede zwischen \textit{er} und \textit{der kleine Junge, der heute Geburtstag hat} nicht behandelt werden. Beide Phrasen können in einem Satz dieselbe Funktion als Subjekt erfüllen: \textit{Jetzt strahlt er} oder eben \textit{Jetzt strahlt der kleine Junge, der heute Geburtstag hat}. Auch die Wortartzugehörigkeit gibt hierüber keinen Aufschluss. Bei \textit{er} handelt es sich um ein Pronomen und bei \textit{der kleine Junge, der heute Geburtstag hat} um die Folge von Artikel, Adjektiv, Nomen, Pronomen, Adverb, Nomen, Verb.

\largerpage
Deshalb fordert \citet[10, 56--60]{Ries1894} neben dem Satzbegriff den Phrasen- bzw. Wortgruppenbegriff zur Beschreibung syntaktischer Strukturen. Angeregt davon überarbeitet \citet{Suetterlin1907} seine erstmals 1897 erschienene Grammatik \gqq{für Lehrer und Studierende}. Dennoch ist mehr als hundert Jahre nach \citet{Ries1894} die Phrase in der Schulgrammatik nicht wirklich angekommen. \citet[10, 14]{Granzow-Emden2014} bezeichnet sie als \gqq{Tor zur Schulgrammatik}. Denn Einsicht in den Aufbau von Phrasen ist sowohl für die Sprachpraxis als auch für die Sprachreflexion relevant: Warum können wir nicht sagen *\textit{ein Apfel roter}, was auf Französisch der Normalfall ist: \textit{une pomme rouge}? Warum heißt es \textit{mit d-em Löffel}, aber \textit{ohne d-en Löffel}, wo es doch eigentlich \textit{d-er Löffel} heißt? Welches Wort steuert den Strukturaufbau in der Wortgruppe \textit{um Bernsteine zu finden} und warum kann es nicht heißen *\textit{um zu finden Bernsteine}, was auf Englisch die einzig mögliche Stellung ist: \textit{to find pieces of amber}?    

\end{Exkurs}



\section{Verbalphrase}\label{section:VP}
\largerpage[2]
Die wesentlichen Grundlagen zum Verständnis der Verbalphrase (VP) haben wir bereits in Kapitel \ref{chapter:Valenz} zum Thema Valenz dargestellt. Das Zentrum der VP ist ein finites Vollverb wie \textit{hoffte} in Satz (\ref{ea:Wiederaufnahme1}).

\ea \label{ea:Wiederaufnahme1} Der Mann hoffte ohne Unterlass auf den Brief.
\z 

Das\is{Verbalphrase (VP)} finite Verb \textit{hoffte} verlangt nach der Anwesenheit anderer Phrasen, nach den Ergänzungen. Der Ausbau des Verbkopfes zur komplexen Phrase geschieht primär durch die Sättigung der Leerstellen. In unserem Satz handelt es sich um die Nominalphrase (NP) \textit{der Mann} und um die Präpositionalphrase (PP) \textit{auf den Brief}. Sekundär wird die Phrase durch Angaben ausgebaut, im Beispiel (\ref{ea:Wiederaufnahme1}) durch die PP \textit{ohne Unterlass}. 

Innerhalb von Phrasen herrschen zwei wesentliche Beziehungen: die in \ref{subsubsection:TestFOSP} besprochene\is{Rektion} \textit{Rektion} als Festlegung von Merkmalen durch den Kopf und die\is{Kongruenz} \textit{Kongruenz}, die Übereinstimmung von Merkmalen zwischen dem Kopf und anderen Teilen der Phrase. 

In Bezug auf den Numerus kongruiert der finite Verbkopf mit der ersten Ergänzung im Nominativ, mit dem Subjekt. Deshalb heißt es \textit{der Mann hoffte} und nicht *\textit{der Mann hofften}. Bei Personalpronomen kongruiert die finite Verbform auch in Bezug auf die Person. Deshalb heißt es \textit{er hoffte} und nicht *\textit{er hofftest}.

Bei Rektion\is{Rektion} geht es primär um die Kasuszuweisung vom Verb an die Ergänzungen. In Bezug auf die erste Ergänzung im Nominativ kann nur von\is{Rektion!kategorial} kategorialer Rektion gesprochen werden, da fast alle finiten aktivischen Verbformen den Nominativ regieren. Die zweite oder dritte Ergänzung ist verbspezifisch regiert. So regiert \textit{hoffte} die Präposition \textit{auf} mit der Forderung nach einer NP im Akkusativ. 

Neben\is{Verbalphrase (VP)} den Ergänzungen kann eine VP durch Phrasen weiter ausgebaut werden, die in \ref{subsection:EvsA} als Angaben klassifiziert wurden. So könnte in (\ref{ea:Wiederaufnahme1}) auch noch \textit{lange} oder \textit{sehr lange} zur Angabe der Dauer eingefügt werden. Die Angabe der Zeitdauer ist nicht durch die Valenz des Verbs festgelegt, gehört jedoch zur VP. 

Die\is{Verbalphrase (VP)} Ergänzungen führen zur Sättigung des Verbkopfes zur VP. Durch Angaben wird die VP erweitert. Ergänzungen und Angaben hängen, wie in Abschnitt~\ref{subsection:EvsA} dargestellt, vom Kopf ab. Ihr Vorkommen ist an das Vorkommen des Kopfes gebunden. 

In den Abschnitten~\ref{section:RegierteVerben} und \ref{section:Verbzweitsätze} haben wir gezeigt, dass Verben im Deutschen häufig komplex sind. Sie bestehen aus einem finiten und einem oder mehreren infiniten oder lexikalischen Bestandteilen. Der Verbkopf kann somit mehrteilig (\textit{analytisch}) sein. Die Verhältnisse innerhalb des Verbkomplexes beschreibt \citet[55, §13]{Tesniere1980deutsch} als \gqq{mikroskopische Vergrößerung} des Ausbaus eines einfachen Verbkopfes zur VP. Denn vom finiten Verb hängen die anderen Verbteile ab, was in Abschnitt~\ref{section:RegierteVerben} dargestellt wurde. Wenn wir im Satz (\ref{ea:Wiederaufnahme1}) den synthetischen Kopf \textit{hoffte} durch den analytischen \textit{hat gehofft} ersetzen, so ist der Verbkopf komplex. Das Finitum \textit{hat} regiert die Form des infiniten Bestandteils \textit{gehofft} zum Ausdruck des Perfekts. Der infinite Bestandteil \textit{gehofft} steuert dem Verbalkomplex seine Valenz und Semantik bei. \citet[55--56]{Eisenberg2020_Satz} bezeichnet das Finitum als grammatischen \textit{Kopf} und die infiniten Teile als semantischen \textit{Kern}.\footnote{Die Unterscheidung zwischen Kern und Kopf ließe sich auch auf synthetische Verbformen wie \textit{hoffte} oder \textit{hofft} übertragen. Der Verbstamm \textit{hoff-} bildet den semantischen Kern und die Flexionsendung \textit{-te} bzw. \textit{-t} den grammatischen Kopf. Das wäre zwar analog zu \textit{gehofft hat}, aber wir würden somit die Wortgrenze unterschreiten und die syntaktische Ebene verlassen.} Syntaktisch komplexe finite Verbköpfe werden in der Abbildung~\ref{fig:komplexerVerbkopf} dargestellt. Dass die finite Verbform der Kopf ist, wird durch Alignierung (senkrechte Verbindungslinie) ersichtlich. Da der Kopf sein Merkmal an den gesamten Komplex weitergibt, trägt der gesamte Komplex das Merkmal \textit{finit} (fin). 

\begin{figure}
\centering
  \begin{forest}
  %l sep+=2em
      [Vfin, calign=last
      [Vinf, tier=preterminal [gehofft \\]]
        [Vinf, tier=preterminal [haben \\
        gehofft\\]]
        [Vfin, tier=preterminal [wird\\
        hat\\, tier=terminal]]] 
  \end{forest}
  \caption{Aufbau des komplexen finiten Verbkopfes}\label{fig:komplexerVerbkopf}
\end{figure}

In Abschnitt~\ref{section:RegierteVerben} haben wir die Rektionsverhältnisse innerhalb von Verbalkomplexen beschrieben. Bei der für Nebensätze typischen Verbletztstellung verläuft die Rektion entgegengesetzt zur linearen Abfolge von rechts nach links.\footnote{Wie in Abschnitt~\ref{section:RegierteVerben} sehen wir hier von der schrittweisen Komplexbildung ab, da es uns nur um die Bestandteile eines komplexen Verbkopfes geht. Aufschluss über die schrittweise Verbindung der Teile des Verbalkomplexes gibt zum Beispiel die Möglichkeit, die infiniten Verbteile als eine Einheit gemeinsam voranzustellen: \textit{Gehofft haben wird Paul lange}. Es verbinden sich die infiniten Verbteile zu einer Einheit, bevor sie mit dem Finitum verbunden werden.}

Neben Rektion weist der Verbkopfkomplex eine weitere syntaktische Beziehung auf, nämlich den\is{Positionsbezug} \textit{Positionsbezug}, den wir im Rahmen des Feldermodells im Kapitel \ref{chapter:Feldermodell} beschrieben haben. In Verbzweit- und Verberstsätzen stehen die Konstituenten eines mehrteiligen Verbkopfes getrennt voneinander, das Finitum in der linken Satzklammer und die infiniten bzw. lexikalischen Teile in der rechten Satzklammer. Im Gegensatz dazu kann der kontinuierliche Aufbau des\is{Verbalphrase (VP)} verbalen Kopfkomplexes und dessen Ausbau durch Ergänzungen und Angaben zur VP bei Verbletztstellung wie in Abbildung~\ref{fig:VPinVerbletzt} durch eine Baumstruktur dargestellt werden. Die Dreiecke stehen als Abkürzungen für komplexe Phrasen, um deren interne Struktur es später geht.

\begin{figure}
\centering
  \begin{forest}
    [VPfin, calign=last
      [NP, tier=preterminal
        [der Mann, narroof]
      ]
      [PP, tier=preterminal
        [ohne Unterlass, narroof]
      ]
      [PP, tier=preterminal
        [auf den Brief, narroof]
      ]
      [Vfin, calign=last
        [Vinf, tier=preterminal [gehofft]]
        [Vfin, tier=preterminal [hat, tier=terminal]]
      ]
    ]
  \end{forest}
  \caption{Aufbau einer finiten VP bei Verbletztstellung}\label{fig:VPinVerbletzt}
\end{figure}

Zuerst verbindet sich das finite Verb \textit{hat} mit dem von ihm regierten infiniten Partizip \textit{gehofft} zu dem komplexen Verbkopf \textit{gehofft hat}. Dieser verbindet sich mit den Präpositionalphrasen \textit{auf den Brief} und \textit{ohne Unterlass} sowie mit der Nominalphrase \textit{der Mann} zu einer finiten Verbalphrase. Bei Verbletztstellung lässt sich die Strukturiertheit kontinuierlich von rechts nach links -- entgegengesetzt zur tatsächlichen Reihenfolge -- beschreiben. Das Ergebnis ist ein Satz, und zwar ein Verbletztsatz, der seinerseits typischerweise von einer nebensatzeinleitenden Subjunktion wie \textit{dass} regiert wird und mit dieser eine Subjunktionalphrase bildet, worum es in Abschnitt~\ref{section:SubjunktP} geht. 

Selbstverständlich lässt sich diese Strukturiertheit auch für Verbzweitsätze darstellen. Dann überschneiden sich wie in der Abbildung~\ref{fig:VPinVerbzweit} die Kanten. Das finite Verb an der zweiten Stelle und der infinite Verbteil am Ende bilden den komplexen, die Satzklammer bildenden Verbkopf. Von diesem hängen die übrigen Phrasen ab und bilden die finite VP. 

\begin{figure}
\centering
\begin{tikzpicture}[every node/.style={align=center}, line width=0.4pt]

  % --- Parameter ---
  \def\leveldist{1.3}    % vertikaler Abstand zwischen Ebenen
  \def\baseoffset{0.34}  % Dreiecke etwas höher über dem Wort
  \def\triapexgap{0.06}  % größerer Abstand zwischen Label und Dreiecksspitze
  \def\baseinset{0.2}    % Basis leicht kürzer als der Text

  % --- Ebene 1: Wörter (kursiv) ---
  \node (A) at (-3.5,0) [anchor=north] {Der Mann};
  \node (B) at (-1.5,0) [anchor=north] {hat};
  \node (C) at (0.5,0)  [anchor=north] {ohne Unterlass};
  \node (D) at (3.0,0)  [anchor=north] {auf den Brief};
  \node (E) at (4.9,0)  [anchor=north] {gehofft};

  % --- Ebene 2: Kategorien ---
  \node (NP)       at ($(A)+(0,\leveldist)$) {NP};
  \node (Vfinword) at ($(B)+(0,\leveldist)$) {Vfin};
  \node (PP1)      at ($(C)+(0,\leveldist)$) {PP};
  \node (PP2)      at ($(D)+(0,\leveldist)$) {PP};
  \node (Vinfword) at ($(E)+(0,\leveldist)$) {Vinf};

  % --- Dreiecke über den Phrasen ---
  \coordinate (A_baseleft)  at ($(A.west)+(\baseinset,\baseoffset)$);
  \coordinate (A_baseright) at ($(A.east)+(-\baseinset,\baseoffset)$);
  \coordinate (A_top)       at ($(NP.south)-(0,\triapexgap)$);
  \draw (A_baseleft) -- (A_top) -- (A_baseright) -- cycle;

  \coordinate (C_baseleft)  at ($(C.west)+(\baseinset,\baseoffset)$);
  \coordinate (C_baseright) at ($(C.east)+(-\baseinset,\baseoffset)$);
  \coordinate (C_top)       at ($(PP1.south)-(0,\triapexgap)$);
  \draw (C_baseleft) -- (C_top) -- (C_baseright) -- cycle;

  \coordinate (D_baseleft)  at ($(D.west)+(\baseinset,\baseoffset)$);
  \coordinate (D_baseright) at ($(D.east)+(-\baseinset,\baseoffset)$);
  \coordinate (D_top)       at ($(PP2.south)-(0,\triapexgap)$);
  \draw (D_baseleft) -- (D_top) -- (D_baseright) -- cycle;

  % --- Verbindungen Wort → Kategorie (leicht gekürzt + luftiger) ---
  \draw ($(B.north)+(0,0.15)$) -- ($(Vfinword.south)-(0,0.15)$);
  \draw ($(E.north)+(0,0.15)$) -- ($(Vinfword.south)-(0,0.15)$);

  % --- Höhere Ebenen ---
  \node (F) at ($(B)+(0,2*\leveldist)$) {Vfin};
  \node (G) at ($(B)+(0,3*\leveldist)$) {VPfin};

  % --- Verbindungen nach oben (alle mit etwas mehr Luft) ---
  \draw ($(Vfinword.north)+(0,0.1)$) -- ($(F.south)-(0,0.1)$);
  \draw ($(Vinfword.north)+(0,0.1)$) -- ($(F.south)-(0,0.1)$);

  \draw ($(NP.north)+(0,0.1)$)  -- ($(G.south)-(0,0.1)$);
  \draw ($(PP1.north)+(0,0.1)$) -- ($(G.south)-(0,0.1)$);
  \draw ($(PP2.north)+(0,0.1)$) -- ($(G.south)-(0,0.1)$);
  \draw ($(F.north)+(0,0.1)$)   -- ($(G.south)-(0,0.1)$);

\end{tikzpicture}
    \caption{Aufbau einer finiten VP bei Verbzweitstellung}\label{fig:VPinVerbzweit}
\end{figure}

Wie wir in Abschnitt~\ref{section:RegierteVerben} gesehen haben, können auch\is{Verbalphrase (VP)!infinit} Infinitive als Köpfe von VPs fungieren. So handelt es sich bei der infiniten VP \textit{ein paar Bernsteine zu finden} in (\ref{ea:wartetmitINFErg}) bei \textit{ein paar Bernsteine} um die zweite Ergänzung von \textit{zu finden}. Im Unterschied zu finiten VPs wird mit infiniten VPs keine Behauptung (Aussage) gemacht. 

\ea \label{ea:wartetmitINFErg} Anna hofft, [ein paar Bernsteine zu finden].
\z 

Die Struktur der\is{Verbalphrase (VP)!infinit} infiniten VP \textit{ein paar Bernsteine zu finden} wird in Abbildung~\ref{fig:infVPinVerbletzt} dargestellt.

\begin{figure} 
  \centering
  \begin{forest}
      [VPinf, calign=last
        [NP, tier=preterminal [ein paar Bernsteine, narroof]]
        [Vinf, tier=preterminal [zu finden]]
]
  \end{forest}
  \caption{Aufbau einer infiniten VP}\label{fig:infVPinVerbletzt}
\end{figure}

Zu den infiniten Verbalphrasen zählen wir auch Partizipien I (\ref{ea:VPinf_Part1}) und II (\ref{ex:VPinf_Part2}) mit ihren Ergänzungen und Angaben, wenn diese nicht wie bei \textit{die} [\textit{ein paar Bernseine findende}] \textit{Anna} oder \textit{die} [\textit{am Strand gefundenen}] \textit{Bernsteine} adjektivierte Bestandteile einer Nominalphrase oder Bestandteile eines analytischen Verbkomplexes wie bei \textit{sie hat ein paar Bernsteine gefunden} sind. 

\ea \label{ea:VPinf_Part}
\ea \label{ea:VPinf_Part1} [Ein paar Bernsteine findend], spaziert Anna am Strand.
\ex \label{ex:VPinf_Part2} [Ein paar Bernsteine gefunden], geht Anna zufrieden nach Hause.
\z 
\z 

Der Aufbau der infiniten Verbalphrasen \textit{ein paar Bernsteine findend} (\ref{ea:VPinf_Part1}) und \textit{ein paar Bernsteine gefunden} (\ref{ex:VPinf_Part2}) sieht genauso aus wie bei der infiniten VP \textit{ein paar Bernsteine zu finden} in Abbildung~\ref{fig:infVPinVerbletzt}.  

Der Grund für unsere Klassifizierung als infinite Verbalphrase liegt darin, dass Partizipien, also \textit{findend} in (\ref{ea:VPinf_Part1}) und \textit{gefunden} in (\ref{ex:VPinf_Part2}) genauso wie Infinitive infinite Verbformen sind. Um die Unterscheidung zwischen finiten und infiniten Verbformen geht es detailliert in Abschnitt~\ref{subsubsection:finit_infinit}. Hier genügt erst einmal, dass eine Verbform dann infinit ist, wenn sie nicht konjugiert ist, d.\,h. wenn sie keine Person-Numerus-Markierung für eine Subjekt-NP hat. Das ist auch der Grund dafür, dass die infiniten Formen des Verbs \textit{finden} ihre erste Ergänzung, den Findenden, nicht an sich binden können. Formal fehlt ihre erste Ergänzung. Wir müssen sie aus dem Kontext erschließen. In allen drei Fällen verstehen wir \textit{Anna} auch als diejenige, die Bernsteine finden will, findet oder gefunden hat.

In grammatischen Modellen, in denen Sätze in Subjekt und Prädikat bzw. Prädikatsverband zerlegt werden, zählt die erste Ergänzung im Nominativ, die Subjekt-NP, nicht zur VP. Begründet wird dies insbesondere mit dem Mangel an verbspezifischer Rektion (vgl. Abschnitt~\ref{subsubsection:TestFOSP}) sowie mit der Tatsache, dass das finite Verb mit dem Subjekt in Person und Numerus kongruiert. In unserem valenziellen Ansatz zählen wir die Subjekt-NP zur Verbalphrase, differenzieren aber zwischen einer finiten VP und einer infiniten VP.   



\section{Nominalphrase}\label{section:NP}

Die einfachsten\is{Nominalphrase (NP)} Nominalphrasen (NP) bilden Eigennamen wie \textit{Fritzi} in (\ref{ea:NPEigennamen}) und Pronomen wie \textit{der} in (\ref{ex:NPPronomen}).  

\ea \label{ea:NPEigennamenPronomen}
\ea \label{ea:NPEigennamen} \textit{Fritzi} bellt immer.
\ex \label{ex:NPPronomen} \textit{Der} bellt immer.
\z 
\z 

Der Eigenname \textit{Fritzi} und das Pronomen \textit{der} verweisen in den Beispielsätzen (\ref{ea:NPEigennamenPronomen}) auf ein bellendes Lebewesen. Eigennamen sind der prototypische Fall von artikellosen\is{Nominalphrase (NP)!artikellos} NPs im Singular. Genauso bilden Pronomen als Kopf alleine eine NP, was in Abbildung~\ref{fig:NPEigennamenPronomen} dargestellt wird. 

\begin{figure}
\centering
\begin{forest}
    [NP
      [N, tier=preterminal
        [Fritzi]
      ]
    ]
\end{forest}
\hspace{5em}
\begin{forest}
    [NP
      [Pro-N, tier=preterminal
        [der]
      ]
    ]
\end{forest}
\caption{\label{fig:NPEigennamenPronomen}Eigennamen und Pronomen als minimale NP}
\end{figure}

Man könnte die durch das Pronomen gebildete Phrase auch Pro-NP nennen, was die Terminologie jedoch unnötig verkomplizieren würde. Denn strukturell sind NPs und Pronomen gleichwertig, weil Pronomen eben an der Stelle von NPs stehen. Während viele Nomen typischerweise mit anderen Wörtern und Phrasen zur NP ausgebaut werden, können Pronomen diese alleine bilden. Sie können weitgehend analog zu den Nomen nach rechts ausgebaut werden -- was wir in diesem Abschnitt betrachten werden -- aber anders als die Nomen nur sehr begrenzt nach links.\footnote{Der Ausbau nach links ist \zB durch sogenannte Fokuspartikeln wie \textit{sogar} (vgl. Abschnitt~\ref{subsubsection:Fokus}) möglich: [\textit{Sogar der}] \textit{bellt.} Nur einige Indefinit- (vgl. Abschnitt~\ref{subsubsection:Indefinitpronomen}) und die Possessivpronomen (vgl. Abschnitt~\ref{subsubsection:PossPronomen}) können mit einem definiten Artikel eine NP bilden: [\textit{Der meine}] \textit{bellt nie}. [\textit{Der eine}] \textit{bellt immer}.}

Artikellos gebrauchte\is{Nominalphrase (NP)!artikellos} Nomen können alleine eine NP bilden. Dies ist standardmäßig beim indefiniten Plural wie \textit{Hund-e} in \textit{Hunde bellen} der Fall. Neben den bereits erwähnten Eigennamen wie \textit{Fritzi} verwenden wir auch sogenannte Massennomen wie \textit{Gold} und Abstrakta wie \textit{Glück} artikellos als NP wie in \textit{Gold glänzt}.   

Typischerweise sind\is{Nominalphrase (NP)} NPs jedoch komplex. Ein Substantiv wie \textit{Hund} bildet als nominaler Bestandteil den Kopf der NP. Dieser Kopf steht in einem bestimmten Numerus und Kasus und er regiert ein Genus. Allerdings sind das Genus und der Kasus nicht primär am Kopf ersichtlich, sondern an einem Artikelwort wie \textit{der} in (\ref{ea:derHund}), \textit{den} in (\ref{ex:denHund}) oder \textit{dem} in (\ref{ex:demHund}).

\ea 
\ea \label{ea:derHund} Jetzt bellt \textit{der Hund} schon wieder.
\ex \label{ex:denHund} Wir hören \textit{den Hund}.
\ex \label{ex:demHund} Wir geben \textit{dem Hund} einen Knochen.
\z 
\z 

Um die verschiedenen Flexionsformen eines Nomens geht es in Abschnitt~\ref{subsection:Substantivflexion}. Jetzt interessiert uns der Aufbau einer minimalen NP aus Artikel und Nomen. Für \textit{der Hund} stellen wir das Schema in der Abbildung~\ref{fig:ArtNNP} dar.\footnote{Es sei erwähnt, dass man in der generativen Grammatik davon ausgeht, dass bei der Verbindung von Artikel und Nomen der Artikel (Determinierer) der Kopf der ganzen Phrase (Determiniererphrase) ist. Eine kurze Übersicht hierzu bietet \citet[96--98]{Tahiri2022} mit entsprechenden Verweisen auf weiterführende Literatur.}  

\begin{figure}  
\centering
  \begin{forest}
      [NP, calign=last
        [Art, tier=preterminal [der]]
        [N, tier=preterminal [Hund]]
]
  \end{forest}
  \caption{Artikel und Substantiv als minimale NP}\label{fig:ArtNNP}
\end{figure}

Erst\is{Artikelfunktion} der Artikel erlaubt es uns, aus einer nominalen Bedeutungsklasse wie \textit{Hund} auf ein Individuum der Gattung oder auf die Gattung selbst zu verweisen.\footnote{\citet[79--82]{Buehler1934} bezeichnet die Bedeutungsklasse als \textit{Symbolfeld} und die Individualisierung als \textit{Zeigefeld} der Sprache.} Auch der indefinite Artikel \textit{ein} in \textit{ein Hund} führt in vielen Fällen zur Individualisierung, indem auf einen beliebigen Vertreter der Klasse verwiesen wird. Artikel grenzen das Substantiv in Bezug auf seinen Geltungsbereich ein.\footnote{\citet[18]{Agel1996} schreibt daher auch vom \textit{finiten Substantiv}. Allerdings macht \citet[32]{Agel1996} deutlich, dass eine echte syntaktische Nominalklammer nicht aus Artikel und Nomen, sondern aus dem Flexionsteil, welcher typischerweise der Artikel trägt, und dem Nomen besteht, also \zB \textit{de}[\textit{r}--\textit{Hund}]. Dies wird deutlich, wenn der Flexionsteil am Adjektiv realisiert wird wie in \textit{klein}[\textit{er}--\textit{Hund}]. Um die Flexion innerhalb der NP geht es in den Abschnitten~\ref{subsection:Artikel_SGr_Flexion} und \ref{subsection:Adjektiv_SGr_Flexion}.}


So wie die Teile des Verbkomplexes in Verberst- und Verbzweitsätzen eine Klammer und somit Felder bilden (vgl. Einführung ins Kapitel \ref{chapter:Feldermodell}), so bilden auch das Artikelwort und das Substantiv eine\is{Nominalklammer} Klammer. Zwischen dem Artikel und dem Nomen können flektierte Adjektive stehen. Die Klammeranalogie wird in der Abbildung~\ref{fig:Nominalklammer} dargestellt. Der linken Satzklammer (LSK) entspricht die linke Nominalklammer (LNK), der rechten Satzklammer (RSK) die rechte Nominalklammer (RNK).   

\begin{figure}
  
  \begin{tabular}{cp{0.1em}cp{0.1em}cp{0.1em}c}
    LSK/LNK && MF && RSK/RNK && \\
    \cmidrule{1-1}\cmidrule{3-3}\cmidrule{5-5}
    hat && der Hund && gebellt \\
    der && kleine && Hund \\
  \end{tabular}
  \caption{Nominalklammer als Analogon zur Verbalklammer}
  \label{fig:Nominalklammer}
\end{figure}


\citet[355--356]{Weinrich2003} beschreibt die\is{Nominalklammer} Nominalklammer als \gqq{Analogon} zur Verbalklammer. So wie das Finitum die Kategorien Person, Numerus, Tempus, Modus und Genus verbi markiert, das Verb damit eingrenzt (Finitum) und das Subjekt via Kongruenz in die Verbgruppe einbindet, so markiert der Artikel die Genus-, Kasus- sowie sekundär Numeruskategorien und stellt typischerweise Referenz auf einen definiten (konkreten) oder indefiniten (beliebigen) Vertreter der vom Substantiv bezeichneten Klasse her. 

\begin{Exkurs}{Funktion der Klammer}
 
\citet[95]{Ronneberger-Sibold2010} beschreibt die grundlegende Funktion der Klammer als \gqq{Erleichterung der syntaktischen Dekodierung} wie folgt: 

\begin{quote}
    
Dadurch, dass die jeweils zueinander passenden Klammerränder die Grenzen von (verschieden definierten) Konstituenten klar markieren, wissen die Hörer/Leser bzw. die Hörerinnen/Leserinnen während des Dekodierungsprozesses jederzeit, ob sie sich am Anfang, im Inneren oder am Ende einer Konstituente befinden. \citep[95]{Ronneberger-Sibold2010}  

\end{quote}
\citet[356]{Weinrich2003} weist auf einen anderen Effekt hin. Der Hörer muss die Informationen aus dem MF solange in seinem Kontextgedächtnis speichern, bis er das Verb in der RSK bzw. das Nomen in der RNK rezipiert. Dies erzeuge einen Spannungsbogen, der durch die rechte Klammer aufgelöst wird. 

Für \citet[223--236]{Bittner2010} haben beide Klammertypen keine eigenständigen Funktionen. Die Klammerbildung resultiere aus der sprachspezifischen Umsetzung von informationsstrukturellen Maximen: die Voranstellung des Finitums kennzeichnet Fragen und Aussagen (vgl. \textit{Illokution} Einführung Kapitel~\ref{chapter:Feldermodell}), die Endstellung des infiniten Vollverbs folgt der Informationsgewichtung, nach der neue Informationen (indefinit bzw. infinit) weit rechts stehen (vgl. Abschnitt~\ref{subsection:Mittelfeld}). Die Satzklammer wird diesen Anforderungen gerecht. Analog zum Finitum eröffnet der Artikel die Nominalklammer und determiniert das rechts stehende Nomen. Damit werden sowohl die Verbal- als auch die Nominalklammer durch ein grammatisches Element eröffnet, das die Bedeutung des am Ende stehenden Nomens aktualisiert.  

Auf die didaktische Nutzung der Nominalklammer bei der Vermittlung der satzinternen Großschreibung (Substantivgroßschreibung) gehen wir gesondert im Exkurs~\ref{ex:Treppenwoerter} in Kapitel~\ref{kap:Graphematik} ein, wo es um die Prinzipien der Schreibung des Deutschen geht.

\end{Exkurs}

Einen Sonderfall in der NP stellen vorangestellte genitivische Eigennamen wie in \textit{Fritzis Fressnapf} und als solche gebrauchte Verwandtschaftsbeziehungen wie in \textit{Omas Liebling} dar. Die Erweiterung durch die pränominale NP im Genitiv hat auch Artikelfunktion. Das ist auch bei den genitivischen Relativpronomen \textit{deren} und \textit{dessen} in \textit{Anna und deren Bruder} oder \textit{Paul und dessen Bruder} der Fall.\footnote{Die Voranstellung des Genitivs ist die historisch ältere Form und geschah vor Ausbildung des Artikels. Bei Ausbildung des Artikelgebrauchs wurde der vorangestellte Genitiv als definites Artikelwort reanalysiert; cf. \citet[231]{Demske2001}. Artikelwort und pränominaler Genitiv können deshalb nicht zusammen vorkommen: *\textit{der Fritzis Fressnapf}. Die vorangestellte genitivische NP besetzt zwar die Stelle des definiten Artikels, ist aber selbst keiner. Das zeigt sich bei der Flexion des Adjektivs (vgl. Abschnitt~\ref{subsection:Adjektiv_SGr_Flexion}): \textit{der neu-e Fressnapf} versus \textit{Fritzis neu-er Fressnapf}. Bei Interesse verweisen wir auf \citet{Eichinger2006}.}

In der\is{Nominalklammer!Mittelfeld} Nominalklammer können flektierte Adjektive wie \textit{kleine} in \textit{der kleine Hund} stehen. Die Adjektive hängen vom Nomen ab. Ihre Funktion ist es, Werte von Eigenschaften festzulegen, über die das Nomen bzw. die von der NP bezeichnete Entität verfügt. So legt \textit{kleine} den Wert der Eigenschaft Größe fest, über die der bezeichnete Hund verfügt. Der Unterschied zur verbalen Prädikation besteht darin, dass mit \textit{der kleine Hund} nicht explizit behauptet wird, dass der Hund klein ist.\footnote{\citet[317]{Eichinger1991} beschreibt die Merkmalsübertragung durch Adjektive als \gqq{implizite Prädikation}. \citet[140]{Paul1880} bezeichnet Merkmalsbestimmungen dieser Art als \gqq{degradiertes Prädikat}.} Sie kann explizit gemacht werden, wenn das Adjektiv mit \textit{sein} prädikativ gebraucht wird wie in \textit{der Hund ist klein}. Darum geht es im Abschnitt~\ref{subsection:SP und OP}. 

Auch Adjektive bilden Phrasen. Wir könnten statt \textit{kleine} in \textit{der kleine Hund} auch \textit{der sehr kleine Hund} sagen. Um den Bau von Adjektivphrasen (APs) geht es in Abschnitt~\ref{section:AP}. Hier interessiert erst einmal nur der Aufbau der NP mit einer AP als Konstituente. Dies wird in Abbildung~\ref{fig:NPmitAP} dargestellt.


\begin{figure}
  \centering
  \begin{forest}
    [NP, calign=child, calign child=3
      [Art, tier=preterminal
        [der]
      ]
      [AP, tier=preterminal
        [sehr kleine, narroof]
      ]
      [N, tier=preterminal
        [Hund]
      ]
      ]
  \end{forest}
  \caption{Ausbau der NP durch eine AP}
  \label{fig:NPmitAP}
\end{figure}

Theoretisch\is{Nominalklammer!Mittelfeld} können beliebig viele APs in der Nominalklammer stehen. Denn man könnte immer wieder eine AP hinzufügen: \textit{der sehr kleine Hund} oder \textit{der freundliche sehr kleine Hund} oder \textit{der erste freundliche sehr kleine Hund} usw.  

\newpage
\begin{Exkurs}{Reihenfolge der Adjektive in der Nominalklammer}
    
\largerpage
Um\is{Nominalklammer!Mittelfeld} die Abfolge der Adjektivphrasen (APs) zu erklären, schlägt \citet[317]{Eichinger1991} ein plurizentrisches Modell vor. In diesem bildet der Artikel einen Pol und das Nomen den anderen. Die APs stehen entweder beim Artikel oder beim Nomen. 

In Artikelnähe stehen Zahladjektive. Denn sie bauen explizit die Quantifikation aus, die im Artikel implizit angelegt ist. Der definite Artikel im Singular wie bei \textit{der Hund} impliziert Einzahl, welche durch ein Zahladjektiv explizit gemacht werden kann wie in \textit{der eine Hund}. Dies ist freilich im Plural viel relevanter, da die Mehrzahl bei \textit{die Hunde} sowohl für \textit{die zwei Hunde} als auch für \textit{die vierhundert Hunde} gelten kann.  

Adjektive\is{Nominalklammer!Mittelfeld}, die eine zeitliche und örtliche Eingrenzung des Bezugsnomens leisten, stehen ebenfalls in Artikelnähe, nach den Zahladjektiven (falls vorhanden). Denn sie vervollständigen die im Artikel angelegte Individualisierung eines konkreten Referenten. Daher heißt es typischerweise \textit{die zwei damaligen Schüler} oder \textit{die zwei hiesigen Schulen}. 

Eine\is{Nominalklammer!Mittelfeld} Art Zwischenstellung zwischen dem Artikelpol und dem Nominalpol nehmen nach \citet[321]{Eichinger1991} die Adjektive ein, die Eigenschaften im Sinne einer impliziten \textit{sein}-Prädikation ausdrücken. Dabei handelt es sich um bewertende Adjektive wie \textit{unkompliziert}, \textit{überraschend} oder \textit{altmodisch} und um Adjektive, die relative Qualitäten wie \textit{groß} oder \textit{alt} ausdrücken. Hierbei handelt es sich um mehr oder weniger typische Eigenschaften von Nominalklassen wie Alter bei Lebewesen oder Größe von Objekten. Wenn beide Adjektivtypen vorkommen, steht das bewertende Adjektiv zuerst, da ja gerade der Komplex aus Eigenschaft und Nomen bewertet wird. Als Beispiele führt \citet[321]{Eichinger1991} Belege wie \textit{diese altmodischen kleinen Gesänge} und \textit{ein entzückendes kleines Palais} an. Denn \textit{die kleinen Gesänge sind altmodisch} und \textit{das kleine Palais ist entzückend}.

Dem nominalen Pol am nächsten stehen\is{Nominalklammer!Mittelfeld} Adjektive, die in der Regel nicht steigerbare Werte absoluter Merkmale festlegen. Bei Konkreta handelt es sich um Merkmale wie Form, Farbe und Material. Deshalb scheint die Abfolge in \textit{ein großer runder Tisch} und \textit{eine kleine blaue Blüte} natürlicher als in \textit{ein runder großer Tisch} und \textit{eine blaue kleine Blüte}. Ebenfalls am nominalen Pol stehen Adjektive, die das Nomen zu Sachbereichen zuordnen wie \textit{die chemische Industrie} oder \textit{der sportliche Ehrgeiz}. Hier wird ein Bereich der Chemie zugeordnet und der Ehrgeiz auf den Sport bezogen. \citet[324]{Eichinger1991} bezeichnet sie als \textit{Nominalklassifikationen}, die eher zu Kategoriebildungen führen.

\end{Exkurs}


Bei\is{Nominalklammer!Mittelfeld} mehreren aufeinanderfolgenden APs wie in (\ref{ea:mehrereAP}) kann es sich trotz formaler Ähnlichkeit um sehr unterschiedliche Strukturen handeln. Man vergleiche (\ref{ea:AP,AP}) mit (\ref{ex:APAP}).  

\ea \label{ea:mehrereAP}
\ea \label{ea:AP,AP} kluge, unterhaltsame Bücher
\ex \label{ex:APAP} spannende französische Bücher
\z 
\z 

Die Bedeutungen von zwei oder mehreren APs können sich wie in (\ref{ea:AP,AP}) gleichrangig (koordiniert) auf ein Nomen beziehen. In diesem Fall setzt man ein Komma. Die koordinierte Lesart kann durch \textit{und} zwischen beiden APs getestet werden (\ref{ea:APgleichrangigmitKOOR}).

\ea \label{ea:APgleichrangigmitKOOR} kluge und unterhaltsame Bücher
\z 

Strukturell\is{Nominalklammer!Mittelfeld} verbinden sich die beiden APs zu einem koordinierten AP-Komplex und dieser wird dann auf das Nomen übertragen. Dies stellen wir in der Abbildung~\ref{fig:koordinierteAPinNP} dar.


\begin{figure}
  \centering
  \begin{forest}
    [NP, calign=last
      [AP, calign=child, calign child=2
        [AP, tier=preterminal
          [kluge\\ kluge, narroof]
        ]
        [Koordination, tier=preterminal
          [{und\\,}]
        ]
        [AP, tier=preterminal
          [unterhaltsame\\unterhaltsame, narroof]
        ]
      ]
      [N, tier=preterminal
        [Bücher\\Bücher]
      ]
    ]
  \end{forest}
  \caption{Koordinierte APs in der NP}
  \label{fig:koordinierteAPinNP}
\end{figure}


Hiervon zu unterscheiden sind Fälle, in denen sich die\is{Nominalklammer!Mittelfeld} APs schrittweise nacheinander auf das Nomen beziehen wie in (\ref{ex:APAP}). Die beiden APs verstehen wir, zumindest in der üblichen Lesart, nicht koordiniert. Zuerst verbindet sich \textit{französische} mit \textit{Bücher} und die AP \textit{spannende} bezieht sich auf den Komplex \textit{französische Bücher.} In dieser üblichen Lesart kann kein \textit{und} zwischen den APs stehen (\ref{ea:APschrittweise*und}). 

\ea \label{ea:APschrittweise*und} ? spannende und französische Bücher
\z 

Die Struktur stellen wir in der Abbildung~\ref{fig:schrittweiseAPAPNP} dar. 

\begin{figure}
  \centering
  \begin{forest}
    %l sep+=2em
    [NP, calign=last
      [AP, tier=preterminal
        [spannende, narroof]
      ]
      [N, calign=last
        [AP, tier=preterminal [französische, narroof]]
        [N, tier=preterminal [Bücher, tier=terminal]]
]
    ]
  \end{forest}
  \caption{Schrittweise Verbindung der APs mit der NP}\label{fig:schrittweiseAPAPNP}
\end{figure}

\begin{Exkurs}{Nachgestellte Adjektive}
    
Zur besonderen Hervorhebung kann das Bezugsnomen alleine vorangestellt werden. Das flektierte Adjektiv steht dann außerhalb der Nominalklammer, also in Distanz zum Kopf wie in (\ref{ea:flektierteAPpostnominal}). Typisch ist ihr Anschluss mit \textit{nur} oder \textit{und zwar}.   

\ea \label{ea:flektierteAPpostnominal} Hier gibt es Bücher ab zwei Euro, und zwar neue.\footnote{Die Presse, 09.06.2008.}
\z 

Auch diese APs sind Konstituenten der\is{Nominalphrase (NP)} NP. Der vorangehende nominale Kopf wird als Bezugsnomen mitverstanden als \textit{und zwar neue Bücher}. 

Ein\is{Nominalphrase (NP)} anderer Sonderfall\is{Nominalklammer!Nachfeld} nachgestellter Adjektive findet sich in dem Kinderlied \textit{Hänschen klein} oder in Goethes \textit{Heidenröslein}, wo es heißt \textit{Röslein rot}. Analoge Bildungen sind \textit{Forelle blau} oder \textit{Apfelkuchen süß und saftig}. Die Adjektive sind nicht flektiert und das Nomen wird artikellos gebraucht. Der Konstruktionstyp gilt als archaisch.\footnote{\citet[29]{Behaghel1927} sagt in einem 1899 gehaltenen Vortrag: \gqq{Wenn das Volkslied noch immer singt von \textit{Röslein rot}, so bewahrt es damit eine Wortstellung, die seit etwa einem Jahrtausend in der lebendigen Sprache abgestorben ist.}} \citet[63]{Duerscheid2002} analysiert diese Stellung funktional: \gqq{Das Element mit dem höheren Grad an kommunikativer Dynamik [das Nomen, MF] steht in der linearen Abfolge an erster Stelle, dann folgt dasjenige mit dem niedrigeren Mitteilungswert, der erläuternde Zusatz}. Die Nicht-Flektiertheit des Adjektivs ist insofern regelkonform, als dass Adjektive als Konstituenten einer NP nur innerhalb der Nominalklammer flektiert werden. Als Ausnahmen hiervon führt die \citet[455, §750]{Duden2022} feste Wendungen aus älteren Sprachstufen auf wie \textit{Gut Ding will Weile haben} und geprägte Verwendungen wie \textit{warm Wasser}. Zu diesem Muster gehören auch \textit{klein Erna} und \textit{lütt Matten}. 

Einen\is{Nominalphrase (NP)} speziellen Typ der unflektierten\is{Nominalklammer!Nachfeld} Nachstellung von Adjektiven zeigen Bildungen wie \textit{Romantik pur}, \textit{Sonne satt}, \textit{Party total} oder \textit{Gesetz light}. Diese nachgestellten Adjektive sind keine Eigenschaftszuweisungen im engeren Sinne, sondern bringen häufig, so \citet[67]{Duerscheid2002}, \gqq{eine wertende Sprechereinstellung zum Ausdruck}. 

Wir kommen unter dem Stichwort \textit{Apposition} im Exkurs~\ref{subsection:Apposition} auf jene Konstruktionen zurück. In Bezug auf die nachgestellten Adjektive gilt, dass die von ihnen gebildeten APs Bestandteile, also Konstituenten von NPs sind, was in \ref{fig:nachgestellteAP} dargestellt wird.

\begin{figure}
  \centering
  \begin{forest}
    [NP, calign=child, calign child=1
      [N, tier=preterminal
        [Hänschen\\Innenwandfarbe\\Romantik]
      ]
      [AP, tier=preterminal
        [klein\\weiß und waschbeständig\\pur]
      ]
      ]
  \end{forest}
  \caption{Ausbau der NP durch eine nachgestellte AP}
  \label{fig:nachgestellteAP}
\end{figure}

\end{Exkurs}

Typischerweise stehen\is{Nominalklammer!Nachfeld} im nominalen Nachfeld, d.\,h. nach dem Kopfnomen, NPs im Genitiv, PPs und Relativsätze. Da innerhalb der Nominalklammer die Flexion der einzelnen Elemente zusammenspielt (\textit{der kleine Mann/ein kleiner Mann}), wird die Ausklammerung von NPs im Genitiv, von PPs und Relativsätzen verständlich. Sie würden die Klammerflexion nämlich blockieren. Das Gleiche wäre der Fall bei Adverbien, die zwar typischerweise Bestandteile von VPs sind, aber auch postnominal wie \textit{der Baum dort} oder \textit{der Film gestern} gebraucht werden. 

Postnominale\is{Nominalklammer!Nachfeld} NPs im Genitiv und PPs ähneln sich funktional, was anhand von \textit{der kleine Hund meiner Schwester} und \textit{der kleine Hund von meiner Schwester} deutlich wird. Relativsätze stehen ganz am Ende der NP.  

Es mag sein, dass die Analogie zum Nachfeld in Sätzen didaktisch hilfreich ist, weshalb wir postnominale Erweiterungen in Abbildung~\ref{fig:postnominalerNPAusbau} darstellen.


\begin{figure}
  
  \begin{tabular}{cp{0.1em}cp{0.1em}cp{0.1em}c}
    LNK && MF && RNK && NF \\
    \cmidrule{1-1}\cmidrule{3-3}\cmidrule{5-5}\cmidrule{7-7}
    der && kleine && Hund && meiner Schwester \\
    der && kleine && Hund && aus Berlin \\
    der && kleine && Hund, && den wir streicheln\\ 
     der && kleine && Hund && meiner Schwester, den wir streicheln\\ 
  \end{tabular}
  \caption{Postnominale Erweiterungen der NP}
  \label{fig:postnominalerNPAusbau}
\end{figure}

In der Abbildung~\ref{fig:komplexeNP} stellen wir die Konstituentenstruktur der\is{Nominalphrase (NP)} NP \textit{der kleine Hund meiner Schwester, den wir streicheln} dar. 


\begin{figure}
  \centering
  \begin{forest}
    [NP, calign=child, calign child=3
      [Art, tier=preterminal
        [der]
      ]
      [AP, tier=preterminal
        [kleine]
      ]
      [N, tier=preterminal
        [Hund]
      ]
      [NP, tier=preterminal
        [meiner Schwester{,}, narroof]
      ]
      [Relativsatz, tier=preterminal
        [den wir streicheln, narroof]
      ]
    ]
  \end{forest}
  \caption{Ausbau der NP durch eine AP, eine NP und einen Relativsatz}
  \label{fig:komplexeNP}
\end{figure}

Grundlegende Unterschiede zwischen dem Ausbau des Verbkopfes zur VP und des Nominalkopfes zur NP sowie deverbale Nominalisierungen behandeln wir im Abschnitt~\ref{section:SatzgliederundAttribute}, in dem es um Satzglieder im Unterschied zu Attributen geht. Bevor wir den Bau von Adjektivphrasen (APs) im Abschnitt~\ref{section:AP} untersuchen, geht es im folgenden Exkurs~\ref{subsection:Apposition} um Appositionen, d.\,h. um komplexe NPs wie \textit{Chaplin, der berühmte Schauspieler} oder \textit{der berühmte Schauspieler Chaplin}.   


\begin{Exkurs}{Appositionen}\label{subsection:Apposition}

Komplexe NPs\is{Apposition} wie \textit{Chaplin, der berühmte Schauspieler} bestehen aus einem Eigennamen, hier \textit{Chaplin}, und einer folgenden NP, hier \textit{der berühmte Schauspieler}. Beide können referenzidentisch sein, d.\,h. sowohl mit \textit{Chaplin} als auch mit \textit{der berühmte Schauspieler} kann dieselbe außersprachliche Person oder Sache bezeichnet werden. Die Funktion dieser Konstruktion ist es, dem vom Eigennamen bezeichneten Individuum charakteristische Merkmale zuzuweisen, um dem Kommunikationspartner, wie \citet[33]{Zifonun2015} schreibt, \gqq{den Anschluss an sein Wissen} zu ermöglichen. 

In (\ref{ea:Apposition}) beziehen sich der Eigenname \textit{Chaplin} und die Erläuterung \textit{der berühmte Schauspieler} auf ein und dieselbe außersprachliche Person. Jene Erläuterungen in Form einer NP nennt man\is{Apposition} \textit{Apposition}. Sie bilden zusammen mit ihrem Bezugswort \textit{eine} NP.\footnote{Welche Strukturen unter \textit{Apposition} fallen und welche nicht, ist stark umstritten und interessiert hier nicht. Eine knappe Übersicht gibt \citet[32--34]{Zifonun2015}.} 

\ea \label{ea:Apposition}
\ea \label{ea:weiteApp} Chaplin, der berühmte Schauspieler, brachte viele zum Lachen.
\ex \label{ex:engeApp} Der berühmte Schauspieler Chaplin brachte viele zum Lachen.
\z 
\z 

Nach \citet[511]{HelbigBuscha} wird meist zwischen lockeren Appositionen wie (\ref{ea:weiteApp}) und\is{Apposition!locker} engen, d.\,h. integrierten Appositionen (\ref{ex:engeApp}) unterschieden. 

In (\ref{ea:weiteApp}) handelt es sich bei \textit{der berühmte Schauspieler} um eine \textit{lockere Apposition}. Sie hebt sich durch eine Sprechpause vom Bezugswort \textit{Chaplin} ab, was schriftlich durch ein Komma gekennzeichnet wird. Die appositive NP gibt zusätzliche Informationen über den Referenten des Bezugsnomens. Sie stellt eine zusätzliche Prädikation (Merkmalszuweisung) dar. 

Typischerweise kongruiert die\is{Apposition!locker} appositive NP hinsichtlich des Kasus des Kopfes.\footnote{\textit{Typischerweise}, da insbesondere beim artikellosen Gebrauch keine Kasusmarkierung möglich ist bzw. der Nominativ als unmarkierter Kasus steht wie in \textit{Chaplin, Schauspieler und Regisseur, brachte viele zum Lachen}. Ebenso gibt es tolerierte Abweichungen wie bei der Angabe des Datums \textit{am Freitag, dem 13.} oder \textit{am Freitag, den 13.} Eventuell wird hier der temporale Akkusativ aus \textit{Berlin, den 13.} kopiert. Wenn das Bezugsnomen im Genitiv steht, wird die appositive NP umgangssprachlich auch im Dativ ausgedrückt wie in dem folgenden Internetbeleg: \textit{die Vita Dieter Müllers, dem Botschafter des außergewöhnlichen Geschmacks}. Laut \citet[444--445, §730]{Duden2022} gilt dieser inkongruente Dativ in der Standardsprache als nicht korrekt. Korrekt müsste es heißen: \textit{die Vita Dieter Müllers, des Botschafters des außergewöhnlichen Geschmacks}.} 
In (\ref{ea:weiteApp}) stehen beide im Nominativ, in (\ref{ea:KongruenzAkk}) im Akkusativ, in (\ref{ex:KongruenzDat}) im Dativ und in (\ref{ex:KongruenzGen}) im Genitiv, wo die Kongruenz durch das Flexiv \textit{-s} sichtbar ist.

\ea \label{ea:KongruenzweiteApposition}
\ea \label{ea:KongruenzAkk} Ohne Chaplin, \textit{den} berühmten Schauspieler, hätten wir weniger gelacht.
\ex \label{ex:KongruenzDat} Mit Chaplin, \textit{dem} berühmten Schauspieler, haben viele gelacht.  
\ex \label{ex:KongruenzGen} Die Autobiographie Chaplin\textit{s}, \textit{des} berühmten Schauspieler\textit{s}, ist weniger bekannt.
\z 
\z 

Die Struktur der\is{Apposition!locker} NPs in (\ref{ea:KongruenzweiteApposition}) stellen wir in der Abbildung~\ref{fig:NPmitAppkongruent} dar. 


\begin{figure}
  \centering
  \begin{forest}
    [NP, calign=child, calign child=1
      [NP, calign=child, calign child=1
        [N, tier=preterminal
        [Chaplin{,}]
      ]
      ]
      [NP, tier=preterminal
        [der berühmte Schauspieler, narroof]
      ]
      ]
  \end{forest}
  \caption{Lockere Apposition}
  \label{fig:NPmitAppkongruent}
\end{figure}

Der Kopf der lockeren Apposition ist der Eigenname \textit{Chaplin}. Eigennamen können, wie wir in Abschnitt~\ref{section:NP} gesehen haben, alleine eine NP bilden. Dies ist auch innerhalb dieser lockeren Apposition der Fall. Der Eigenname ist erweiterbar: \textit{der kleine Chaplin, der berühmte Schauspieler}. Auch die Apposition \textit{der berühmte Schauspieler} ist eine NP. Jede der beiden NPs kann anstelle der anderen NP gebraucht werden: \textit{Chaplin brachte viele zum Lachen} oder eben \textit{der berühmte Schauspieler brachte viele zum Lachen}. Allerdings sind beide NPs nicht gleichrangig wie in einer koordinierten Struktur. Denn bei Koordination durch \textit{und} wird die Referenzidentität aufgehoben: \textit{Chaplin und der berühmte Schauspieler}. Die appositive NP ist der ersten NP locker beigeordnet. 

Wir folgen \citet[2037]{ZifHoffStr1997} und zählen auch nachgestellte Adjektive bzw. Adjektivphrasen wie in (\ref{ea:unflAPalsAPP}) zu den\is{Apposition} Appositionen. Auch bei diesen handelt es sich um eine Art Nachtrag zum nominalen Kopf der gesamten NP. Bildlich beschreibt \citet[777]{Agel2017} diese als \gqq{Verlängerung} der NP und \gqq{Wortgruppenrandglied}. \citet[531]{Weinrich2003} hält jenen Nachtrag oder Anbau (statt Integration) dadurch für motiviert, dass die Nominalklammer entlastet wird und dass jene Nachträge \gqq{in ihrer Bedeutung auffällig gemacht} werden. Typischerweise handelt es sich um zwei koordinierte Adjektive wie in (\ref{ea:unflAPalsAPP}) oder ein erweitertes Partizip wie in (\ref{ex:unflPart}).

\ea \label{ea:nachgestellte AP}
\ea \label{ea:unflAPalsAPP} Wir bieten eine komfortable Wohnung, hell und ruhig, zur Zwischenmiete.
\ex \label{ex:unflPart} Bieten tolle 3-Zi-Wohnung, frisch renoviert, mit Balkon und EBK.
\z
\z 

Lockere Appositionen\is{Apposition!locker} sind nur schwach in die NP des Bezugsnomens integriert. Dies wird deutlich, wenn sie nicht auf das Bezugsnomen folgen, sondern von ihm getrennt im Nachfeld eines Satzes erscheinen (\ref{ea:APPimNFbeiNomen}). Typisch ist die Stellung im Nachfeld, wenn das Bezugswort ein Pronomen ist (\ref{ex:APPimNFbeiPrn}). Genauso können auch APs nachgestellt werden (\ref{ex:APPimNFmitAP}).

\ea \label{ea:APPimNF}
\ea \label{ea:APPimNFbeiNomen} Chaplin brachte viele zum Lachen, dieser berühmte Schauspieler.
\ex \label{ex:APPimNFbeiPrn} Er brachte viele zum Lachen, der berühmte Schauspieler.
\ex \label{ex:APPimNFmitAP} Wir bieten eine komfortable Wohnung zur Zwischenmiete, hell und ruhig. 
\z 
\z 

Sogenannte \textit{enge} Appositionen\is{Apposition!eng} wie \textit{der Schauspieler Chaplin} verhalten sich anders. Hier ist \textit{Schauspieler} der Kopf, in den der appositive Eigenname integriert ist. Dieser ist kein Nachtrag oder Zusatz, sondern referenznotwendig. Ohne ihn wüssten wir nicht, welche Person oder Entität bezeichnet wird. Deshalb bezeichnet die \citet[448, §737]{Duden2022} die Funktion dieser engen Apposition als \textit{explikativ}. 

Die Kasusmarkierung der gesamten NP zeigt sich am Artikel und im Genitiv auch am Kopfnomen: \textit{des Schauspielers}. Der appositive Eigenname bildet hier selbst keine Phrase. Es handelt sich um ein nicht erweiterbares Nomen: *\textit{der berühmte Schauspieler kleiner Chaplin}. Der Eigenname \textit{Chaplin} bleibt unverändert im Nominativ als unmarkierter Kasus. Beim Genitiv wird dies beim Vergleich von \textit{Chaplins} in (\ref{ex:KongruenzGen}), hier als (\ref{ea:KongruenzGen1}) wiederholt, mit \textit{Chaplin} in (\ref{ex:keineKongruenzGen}) deutlich.

\ea 
\ea \label{ea:KongruenzGen1} Die Autobiographie \textit{Chaplins}, des berühmten Schauspielers, ist weniger bekannt. 
\ex \label{ex:keineKongruenzGen} Die Autobiographie des berühmten Schauspielers \textit{Chaplin} ist weniger bekannt.
\z 
\z 

\citet[281]{Eisenberg2020_Satz} bezeichnet die Nominativrektion als konstruktionell. Denn der Nominativ wird weder von der Kategorie Nomen noch von einzelnen Nomen regiert, sondern von der Appositionskonstruktion.\footnote{Kategorial regieren Nomen nur den Genitiv wie in \textit{der Hund des Nachbarn}.} Das Gleiche ist der Fall bei vorangestellten Verwandtschafts-, Titel- oder Berufsbezeichnungen, wenn diese einen Artikel haben wie: \textit{mein Onkel Paul}, \textit{der bekannte Professor Eisenberg} oder \textit{der beliebte Bäcker Hacker}. Nur der Artikel und im Genitiv auch das Kopfnomen tragen die Kasusmarkierung der gesamten NP (\ref{ea:meinesOnkelsPaul}). Der appositive Eigenname muss im Nominativ stehen, weshalb (\ref{ex:meinesOnkels*Pauls}) ungrammatisch ist.

\ea \label{ea:meinOnkelPaul}
\ea [] {das Haus meines Onkels \textit{Paul}}\label{ea:meinesOnkelsPaul} 
\ex [*] {das Haus meines Onkels \textit{Pauls}}\label{ex:meinesOnkels*Pauls}
\z 
\z 

Den\is{Apposition!eng} Aufbau der NPs \textit{der Schauspieler Chaplin} und \textit{mein Onkel Paul} stellen wir in Abbildung~\ref{fig:NPmitAppNominativ} dar.


\begin{figure}
\centering
  \begin{forest}
    [NP, calign=child, calign child=2
      [Art, tier=preterminal
        [der\\mein]
      ]
      [N, tier=preterminal
        [Schauspieler\\ Onkel]
      ]
      [N, tier=preterminal
        [Chaplin\\Paul]
      ]
      ]
  \end{forest}
  \caption{Enge Apposition}
  \label{fig:NPmitAppNominativ}
\end{figure}


Diese\is{Apposition!eng} Verhältnisse drehen sich um, wenn das erste Nomen zum Ausdruck einer Verwandtschafts-, Titel- oder Berufsbezeichnung artikellos gebraucht wird, wie in \textit{Onkel Paul}, \textit{Professor Eisenberg} oder \textit{Bäcker Hacker}. In diesen Fällen ist der Eigenname der Kopf der NP, denn er trägt die Kasusmarkierung der NP. Der Nominativ des appositiven Nomens ist konstruktionell regiert. Dies wird wieder beim Genitiv deutlich (\ref{ea:OnkelPaul}). Der Eigenname als Kopf trägt die Genitivmarkierung (\ref{ea:OnkelPauls}). Der appositive Eigenname muss im Nominativ stehen, weshalb (\ref{ex:Onkels*Pauls}) ungrammatisch ist.

\ea \label{ea:OnkelPaul}
\ea [] {das Haus \textit{Onkel} \textit{Pauls}}\label{ea:OnkelPauls}
\ex [*] {das Haus \textit{Onkels} \textit{Pauls}}\label{ex:Onkels*Pauls}
\z 
\z  

Den Aufbau der NPs \textit{Onkel Paul} und \textit{Professor Eisenberg} zeigen wir in Abbildung~\ref{fig:APPmitEigennamenKopf}. Da sich die Verbindung von Vor- und Nachnamen wie \textit{Peter Eisenberg} genauso gestaltet, führen wir sie mit auf.


\begin{figure}
  \centering
  \begin{forest}
    [NP, calign=child, calign child=2
      [N, tier=preterminal
        [Onkel\\ Professor\\ Peter]
      ]
      [ N, tier=preterminal
        [Paul\\ Eisenberg\\ Eisenberg]
      ]
      ]
  \end{forest}
  \caption{Enge Apposition}
  \label{fig:APPmitEigennamenKopf}
\end{figure}

\newpage
In \textit{Onkel Paul} ist \textit{Paul} der Kopf und \textit{Onkel} die\is{Apposition!eng} Apposition. Im Gegensatz dazu ist in \textit{mein Onkel Paul} der Eigenname \textit{Paul} die Apposition und \textit{Onkel} der Kopf. Verständlich wird jenes \gqq{Umkippen der Abhängigkeitsverhältnisse}, wie \citet[282]{Eisenberg2020_Satz} schreibt, vor dem Hintergrund der analytischen Nominalflexion. Wie wir gesehen haben, sind minimale NPs typischerweise zweigliedrig, sie bestehen aus einem Artikel und dem Nomen wie \textit{mein Onkel} und \textit{der Professor}. Wenn der Artikel jedoch fehlt wie in \textit{Onkel Paul}, dann kann das Nomen \textit{Onkel} im Genitiv nicht mehr das Nominalflexiv \textit{-s} tragen, da dessen Realisierung an die analytische Flexion mit Artikel gebunden ist. Daher heißt es \textit{das Haus mein-es Onkel-s Paul} aber nicht *\textit{das Haus Onkel-s Paul}. In diesem Fall trägt der Eigenname das Genitivflexiv und ist Kopf der NP: \textit{das Haus Onkel Paul-s}. Denn Eigennamen verfügen unabhängig von ihrem Genus über synthetische Genitivformen wie \textit{Pauls} und \textit{Annas}.\footnote{\citet{Agel1996} gibt eine knappe Übersicht zum Wandel der synthetischen zur analytischen Substantivflexion.}   

Um enge\is{Apposition!eng} Appositionen handelt es sich auch bei komplexen NPs wie \textit{zwei Flaschen Wasser}. Ihre erste Konstituente besteht aus einem Zahladjektiv und einer Maß-, Mengen- oder Behälterbezeichnung. Ihr Zweitglied ist eine Stoff- oder Inhaltsangabe, die durch die Konstruktion gemessen wird.\footnote{\citet[515]{HelbigBuscha} bezeichnen solche Konstruktionen als \gqq{appositionsverdächtig}. \citet[282--283]{Eisenberg2020_Satz} beschreibt sie unter \textit{engen Appositionen}. Allerdings bemerkt \citet[280]{Eisenberg2020_Satz} vorab, dass es ihm nicht um die Klärung des umstrittenen Appositionsbegriffs gehe, sondern um die Beschreibung von Strukturen, die häufig als Apposition bezeichnet werden.} 

Vorläufer dieser\is{Apposition!eng} engen Apposition ist der partitive Genitiv zum Ausdruck von Teil-Ganzes-Beziehungen\footnote{Der Begriff \textit{partitiv} ist vom lateinischen Wort \textit{pars} (Teil) abgeleitet. Im Deutschen können partitive Beziehungen durch den Genitiv erfolgen wie in \textit{die Hälfte des Wassers} oder durch die Präposition \textit{von} wie in \textit{die Hälfte vom Wasser}.} wie in (dem heute nicht mehr akzeptablen Ausdruck) *\textit{eine Flasche Wassers} statt \textit{eine Flasche Wasser}. \citet[294]{Paul1919} weist auf Belege aus dem Neuhochdeutschen wie \textit{ein Stück Ackers} oder \textit{zehn Pfund Silbers} hin. Aufgrund der analytischen Nominalflexion ist dies heute nicht mehr möglich bzw. gebräuchlich. Die Flexion am Nomen ist nämlich an die Flexion am Artikel oder Adjektiv gebunden. Daher ginge \textit{eine Flasche frischen Wassers} aber nicht mehr *\textit{eine Flasche Wassers} (vgl. Abschnitt~\ref{subsubsection:st_sw_SFL} und \ref{subsection:Artikel_SGr_Flexion}). 

Für die Entwicklung von \textit{eine Flasche Wassers} zur heute üblichen engen Apposition \textit{eine Flasche Wasser} macht \citet[294--295]{Paul1919} Substantive verantwortlich, deren Genitivform sich nicht vom Nominativ unterscheidet. Dies ist der Fall bei Feminina im Singular wie bei \textit{ein Löffel Suppe} oder \textit{ein Pfund Gerste} sowie bei allen drei Genera, wenn der Plural artikellos und ohne vorangestellte AP gebraucht wird wie in \textit{ein Pfund Äpfel}.\footnote{Sobald ein Adjektiv die Inhaltsangabe näher bestimmt, kann durch die Flexion am Adjektiv auch wieder der gehoben klingende Genitiv gebraucht werden: \textit{ein Löffel heißer Suppe} statt des üblicheren \textit{ein Löffel heiße Suppe}, \textit{ein Pfund roter Äpfel} statt als Apposition \textit{ein Pfund rote Äpfel}.} Hiervon ausgehend wurde die Inhaltsangabe nicht mehr als Genitiv-, sondern als neutrale Grundform empfunden. 

Für partitive NPs im Genitiv ist die Analyse klar: der nominale Kopf ist die Maßangabe wie \textit{Flasche} und von dieser abhängig ist die NP im Genitiv wie \textit{frischen Wassers} in \textit{eine Flasche frischen Wassers}. Dieser partitive Genitiv wird, wenn überhaupt, nur noch nach Sammelbezeichnungen gebraucht, wenn das appositive Substantiv im Plural steht und zusammen mit einem Artikel oder mit einem flektierten Adjektiv eine NP bildet: \textit{eine Menge kleiner Probleme} oder \textit{eine Gruppe junger Leute}. Bei solchen Konstruktionen wird anstelle der appositiven NP im Genitiv auch eine appositive PP mit \textit{von} oder \textit{an} gebraucht: \textit{eine Menge von}/\textit{an kleinen Problemen} oder \textit{eine Gruppe von jungen Leuten}.

Der heute gängige Fall ist die enge Apposition wie bei \textit{eine Flasche Wasser}. Das appositive Nomen \textit{Wasser} bildet eine NP. Ersichtlich wird das beim Ausbau durch Adjektive wie in \textit{eine Flasche frisches Wasser}. Die starke Flexion des Adjektivs \textit{frisches} entspricht dem bestimmten Artikel \textit{das}: \textit{das frische Wasser} aber \textit{frisches Wasser}. Um die Adjektivflexion geht es in Abschnitt~\ref{subsection:Adjektiv_SGr_Flexion}.

Bei der heute gängigen engen partitiven Apposition kongruiert die appositive NP mit dem Kasus des Kopfnomens. Das wird sichtbar, wenn das Kopfnomen und das Nomen der appositiven NP maskulin sind. Deshalb zeigen wir die Kongruenz anhand der komplexen NP \textit{ein Krug frischer Saft} im Nominativ (\ref{ea:KrugSaftNom}), im Akkusativ (\ref{ex:KrugSaftAkk}), im Dativ (\ref{ex:KrugSaftDat}) und im Genitiv (\ref{ex:KrugSaftGen}).\footnote{\citet[284]{EisenbergSayatz2002} weist darauf hin, dass auch inkongruente Formulierungen wie \textit{mit einem Krug frischer Saft} oder \textit{mit einem Krug frischen Saft} für manche Sprecherinnen und Sprecher akzeptabel seien.}

\ea 
\ea \label{ea:KrugSaftNom} [Ein Krug frischer Saft] steht auf dem Tisch.
\ex \label{ex:KrugSaftAkk} Wir bestellen [einen Krug frischen Saft].
\ex \label{ex:KrugSaftDat} Nach [einem Krug frischem Saft] hat man keinen Durst mehr.
\ex \label{ex:KrugSaftGen} Anstatt [eines Krugs frischen Safts] trinken sie lieber Wasser. 
\z 
\z 

Den Aufbau dieser komplexen NPs zeigen wir in Abbildung~\ref{fig:engeAppPart}.

\begin{figure}
  \centering
  \begin{forest}
    [NP, calign=child, calign child=2
      [Art, tier=preterminal
        [ein]
      ]
      [N, tier=preterminal
        [Krug]
      ]
      [NP, tier=preterminal
        [frischer Saft, narroof]
      ]
      ]
  \end{forest}
  \caption{Enge partitive Apposition (Mengen- und Maßangabe)}
  \label{fig:engeAppPart}
\end{figure}

Laut \citet[446, §734]{Duden2022} gelten die folgenden Abweichungen von der Kongruenzregel ebenfalls als korrekt. Wenn die gesamte NP im Dativ steht und das Nomen der appositiven NP im Plural, dann kann, der Regel folgend, Kongruenz wie in (\ref{ea:KorpÄpfelDat}) bestehen. Die appositive NP kann aber auch neutral im Nominativ stehen wie in (\ref{ex:KorbÄpfelNom}). Parallel dazu gilt die ältere Konstruktionsweise mit der appositiven NP im Genitiv wie in (\ref{ex:KorbÄpfelGen}) als richtig.

\ea 
\ea \label{ea:KorpÄpfelDat} Paul kam mit [einem Korb frischen Äpfeln].
\ex \label{ex:KorbÄpfelNom} Paul kam mit [einem Korb frische Äpfel].
\ex \label{ex:KorbÄpfelGen} Paul kam mit [einem Korb frischer Äpfel].
\z 
\z 
 
Diese Variation ändert aber nichts am Aufbau der gesamten Phrase mit \textit{Korb} als Kopf. Dessen Kopfstatus zeigt sich auch in der Numeruskongruenz zwischen dem finiten Verb und der gesamten NP \textit{ein Korb Äpfel} als dessen Subjekt. Die gesamte NP \textit{ein Korb frischer Äpfel} steht nach dem Kopfnomen \textit{Korb} im Singular und nicht nach dem Plural von \textit{Äpfel} im Plural. Deshalb steht auch das Finitum wie in (\ref{ea:MaßAPPVerbkongruenz}) im Singular und nicht wie in (\ref{ex:InhaltAPPVerbnichtkongruent}) im Plural.

\ea \label{ea:MengenAPPVerbkongruenz}
\ea [] {Ein Korb Äpfel steht auf dem Tisch.}\label{ea:MaßAPPVerbkongruenz}
\ex [*] {Ein Korb Äpfel stehen auf dem Tisch.}\label{ex:InhaltAPPVerbnichtkongruent}
\z 
\z 

Allerdings wird\is{construction ad sensum} im alltäglichen Sprachgebrauch hiervon abgewichen. Das semantische Gewicht der Inhaltsangabe führt dazu, dass Sprecherinnen und Sprecher Numeruskongruenz nicht mit dem Kopf, sondern mit der NP herstellen, die den Inhalt, das Gemessene ausdrückt. Die \citet[136--137, §124]{Duden2022} führt Beispiele auf wie \textit{drei Kilo Brot reicht aus} statt \textit{drei Kilo Brot reichen aus}. Im ersten Fall wird gegen eine grammatische Regel verstoßen und nach dem empfundenen Sinn konstruiert. Die NP \textit{drei Kilo Brot} wird so verstanden, dass es primär um Brot geht und nicht um die Mengenangabe. Diesem Verständnis nach wird Kongruenz zwischen \textit{Brot} und \textit{reicht} hergestellt, statt zwischen \textit{drei Kilo} und \textit{reichen}. Solche Regelabweichungen werden auch als \textit{constructio ad sensum} bezeichnet. 

\end{Exkurs}


\section{Adjektivphrase}\label{section:AP}

Adjektive können alleine wie in \textit{der} [\textit{kleine}] \textit{Mann} Adjektivphrasen\is{Adjektivphrase (AP)} (AP) bilden, was wir in Abbildung~\ref{fig:minimaleAP} darstellen.

\begin{figure}
\centering
\begin{forest}
    [AP
      [A, tier=preterminal
        [kleine]
      ]
    ]
\end{forest}
 \caption{Adjektiv \textit{kleine} als minimale AP}
  \label{fig:minimaleAP}
\end{figure}


Adjektive können zu Wortgruppen ausgebaut werden, \zB durch Steigerungspartikeln wie \textit{sehr} in \textit{der} [\textit{sehr kleine}] \textit{Mann} oder \textit{ziemlich} in \textit{der} [\textit{ziemlich kleine}] \textit{Mann}. Analog dazu sind unflektierte Adjektive wie \textit{unglaublich} in [\textit{unglaublich schöne}] \textit{Blumen} zu analysieren. Sie bauen den adjektivischen Kopf \textit{schöne} aus. Die Konstituentenstruktur stellen wir in Abbildung~\ref{fig:APmitPartikel} dar. 

\begin{figure}
  \centering
  \begin{forest}
    [AP, calign=last
      [Partikel, tier=preterminal
        [sehr\\
        unglaublich]
      ]
      [A, tier=preterminal
        [kleine\\
        schöne]
      ]
    ]
  \end{forest}
  \caption{AP mit Steigerungspartikel}
  \label{fig:APmitPartikel}
\end{figure}

Einen Sonderfall stellen Erweiterungen der AP wie in (\ref{ea:AdjErgKasusPräp}) dar, die in der\is{Adjektivvalenz} Valenz des jeweiligen Adjektivs angelegt sind. 

\ea \label{ea:AdjErgKasusPräp}
\ea \label{ea:AdjErgKasus} der [seinem Herrchen treue] Hund
\ex \label{eX:AdjErgPP} das [auf seinen Hund stolze] Herrchen
\ex \label{ex:AdjGenitiv} ein [keines klaren Gedankens fähiger] Mensch
\ex \label{ex:AdjGenPräp} ein [zu keinem klaren Gedanken fähiger] Mensch
\z 
\z 

Analog zur Verbrektion weist der adjektivische Kopf seiner Ergänzung entweder einen Kasus wie den Dativ in (\ref{ea:AdjErgKasus}) oder eine Präposition wie \textit{auf} in (\ref{eX:AdjErgPP}) zu. Adjektive können auch wie in (\ref{ex:AdjGenitiv}) den Genitiv regieren, der jedoch wie in (\ref{ex:AdjGenPräp}) einer Präposition weicht.

Den Aufbau von\is{Adjektivphrase (AP)} APs mit Ergänzungen stellen wir anhand von \textit{auf seinen Hund stolze} aus (\ref{eX:AdjErgPP}) in Abbildung~\ref{fig:APmitErg} dar. 

\begin{figure}
  \centering
  \begin{forest}
    [AP, calign=last
      [PP, tier=preterminal
        [auf seinen Hund, narroof]
      ]
      [A, tier=preterminal
        [stolze]
      ]
    ]
  \end{forest}
  \caption{Ergänzungen in der AP}
  \label{fig:APmitErg}
\end{figure}

Ähnlich wie die VP kann auch die\is{Adjektivphrase (AP)} AP über die Ergänzungen hinaus schrittweise weiter ausgebaut werden. Es handelt sich um Angaben \zB zur Zeit, zum Ort oder zum Grund usw., die die vom Adjektiv ausgedrückte Eigenschaft näher bestimmen. So lässt sich die AP in (\ref{eX:AdjErgPP}) problemlos durch Präpositional- und Adverbphrasen sowie eine Steigerungspartikel wie in (\ref{ea:APmitErgundAng}) ausbauen.

\ea \label{ea:APmitErgundAng} das [heute aus gutem Grund auf seinen Hund sehr stolze] Herrchen
\z 

Analog zur Darstellung in Abbildung~\ref{fig:APmitErg} bauen auch die nicht valenziell im Adjektivkopf angelegten Phrasen die\is{Adjektivphrase (AP)} AP aus, was wir in Abbildung~\ref{fig:komplAP} darstellen. Der adjektivische Kopf regiert (wie Verben) nach links. 

\begin{figure}
  \centering
  \begin{forest}
    [AP, calign=last
      [AdvP, tier=preterminal
        [heute]
        ] 
      [PP, tier=preterminal
        [aus gutem Grund, narroof]
      ]
      [PP, tier=preterminal
        [auf seinen Hund, narroof]
      ]
      [Partikel, tier=preterminal
        [sehr]
      ]
      [A, tier=preterminal
        [stolze]
      ]
    ]
  \end{forest}
  \caption{Komplexe AP}
  \label{fig:komplAP}
\end{figure}

Auf adjektivisch gebrauchte Partizipien wie \textit{der auf den Brief wartende Mann} oder \textit{das von Paul gebaute Haus} gehen wir in Abschnitt~\ref{section:SatzgliederundAttribute} genauer ein. 


\section{Präpositionalphrase}\label{section:PP}

Präpositionalphrasen\is{Präpositionalphrase (PP)} (PP) sind Wortgruppen, deren Kopf eine Präposition ist. Um Präpositionen als Wortart geht es im Abschnitt~\ref{subsection:Präposition}. Hier interessiert uns der Bau von Präpositionalphrasen.  

Präpositionen verlangen typischerweise rechts von sich nach einer NP im Akkusativ, im Dativ oder im Genitiv. Deshalb heißt es beispielsweise \textit{mit der Bahn} und nicht *\textit{mit die Bahn}. Die Struktur der PP \textit{mit der Bahn} stellen wir in Abbildung~\ref{fig:PPAufbau} dar.  

\begin{figure}
  \centering
  \begin{forest}
    [PP, calign=child, calign child=1
      [P, tier=preterminal
        [mit]
      ]
      [ NP, tier=preterminal
        [der Bahn, narroof]
      ]
      ]
  \end{forest}
  \caption{Aufbau einer typischen PP}
  \label{fig:PPAufbau}
\end{figure}

In Abschnitt~\ref{section:NP} haben wir gesehen, dass NPs sehr komplex sein können. Das gilt auch für NPs, die Teil einer\is{Präpositionalphrase (PP)} PP sind wie \textit{mit der historischen Bahn aus dem 19. Jahrhundert, die noch mit Dampf betrieben wird}. Das ändert nichts am Bau der PP, was in Abbildung~\ref{fig:PPmitkomplexerNP} gezeigt wird.

\begin{figure}
  \centering
  \begin{forest}
    [PP, calign=child, calign child=1
      [P, tier=preterminal
        [mit]
      ]
      [ NP, tier=preterminal
        [der historischen Bahn\\aus dem 19. Jahrhundert{,}\\die noch mit Dampf betrieben wird, narroof]
      ]
      ]
  \end{forest}
  \caption{Aufbau einer PP mit komplexer NP}
  \label{fig:PPmitkomplexerNP}
\end{figure}

Wenn die von der\is{Präpositionalphrase (PP)!Pro-Form} Präposition regierte NP wie in (\ref{ea:anseineOma}) auf ein Lebewesen verweist, so kann sie wie in (\ref{ex:anihn}) durch ein Pronomen realisiert werden. Die Struktur der PP bleibt unverändert, da wir Pronomen im Abschnitt~\ref{section:NP} syntaktisch als NP analysiert haben.   

\ea \label{ea:PPPronomen}
\ea \label{ea:anseineOma} Anna erinnert sich [an ihren Opa].
\ex \label{ex:anihn} Anna erinnert sich [an ihn].
\z
\z 

Wenn\is{Präpositionalphrase (PP)!Pro-Form} die von der Präposition regierte NP nicht auf ein Lebewesen verweist und die Präposition von der Verbvalenz gefordert wird, so wird bei besonders häufig gebrauchten Präpositionen die gesamte PP durch eine Pro-Form mit \textit{da}(\textit{r})- bzw. bei Fragen und in Relativsätzen mit \textit{wo}(\textit{r})- ersetzt: \textit{sich woran/daran erinnern}.\footnote{Dies ist bei den Präpositionen mit Akkusativ \textit{durch}, \textit{für}, \textit{gegen}, \textit{um}, bei den Präpositionen mit Dativ \textit{aus}, \textit{bei}, \textit{mit}, \textit{nach}, \textit{von}, \textit{zu} und bei den sogenannten Wechselpräpositionen möglich, um die es im Abschnitt~\ref{subsection:Präposition} geht. Alle anderen Präpositionen bleiben unverändert und die NP wird pronominalisiert.} Das \textit{-r-} wird eingeschoben, wenn die Präposition vokalisch beginnt, um das Aufeinandertreffen zweier Vokale, die zu verschiedenen Silben gehören, zu vermeiden. So wird die PP in (\ref{ea:andasKonzert}) durch die Pro-Form \textit{daran} in (\ref{ex:daran}) ersetzt.\footnote{Wenn das Nomen ein Nicht-Lebewesen bezeichnet und ein Neutrum im Dativ (\ref{ea:Neutr.Dativ}) ist oder ein Maskulinum bzw. Femininum im Dativ (\ref{ex:Mask.Dativ}, \ref{ex:Fem.Dativ}) oder im Akkusativ (\ref{ex:Mask.Akk}, \ref{ex:Fem.Akk}), dann wird alternativ die\is{Präpositionalphrase (PP)!Pro-Form} Pro-Form genutzt oder aber die NP als Pronomen realisiert.

\ea \label{ea:ProPräpoderPronomen} 
\ea \label{ea:Neutr.Dativ} Anna zweifelt [an dem Konzept]/[daran]/[an ihm].
\ex \label{ex:Mask.Dativ} Anna hängt [an ihrem Garten]/[daran]/[an ihm].
\ex \label{ex:Fem.Dativ} Anna hängt [an dieser Erinnerung]/[daran]/[an ihr].
\ex \label{ex:Mask.Akk} Anna erinnert sich [an den ersten Schultag]/[daran]/[an ihn].
\ex \label{ex:Fem.Akk} Anna erinnert sich [an ihre Kindheit]/[daran]/[an sie].
\z 
\z  
}

\ea \label{ea:PPPronomenUnbelebt}
\ea \label{ea:andasKonzert} Anna erinnert sich [an das Konzert].
\ex \label{ex:daran} Anna erinnert sich [daran].
\z
\z 

Die\is{Präpositionalphrase (PP)!Pro-Form} Pro-Form bildet allein und direkt die PP, was in Abbildung~\ref{fig:minimalePP} gezeigt wird. 

\begin{figure}
\centering
\begin{forest}
    [PP
        [daran]
    ]
\end{forest}
 \caption{Pro-Form als minimale PP}
  \label{fig:minimalePP}
\end{figure}

Weniger typisch sind im Deutschen sogenannte\is{Postposition} \textit{Postpositionen}. Sie fordern als Köpfe links von sich eine NP wie \textit{der Ordnung halber}. Manche Präpositionen werden auch als Postposition gebraucht, \zB \textit{nach meiner Meinung} oder \textit{meiner Meinung nach}, \textit{des schönen Wetters wegen} oder \textit{wegen des schönen Wetters}. An den Rektionsverhältnissen ändert sich nichts. Der Kasus der NP wird von der Postposition regiert.

Selten sind im Deutschen Komplexe aus Präposition und Postposition, deren NP in der Mitte steht wie \textit{um des lieben Friedens willen}, \textit{von Montag an} oder \textit{von Berlin aus}. Diese Komplexe bezeichnet man auch als \textit{Zirkumposition}.  


Einige Präpositionen nehmen außer NPs auch Adjektive bzw. APs wie \textit{etwas} [\textit{für gut}] \textit{befinden}, Adverbien bzw. Adverbphrasen wie \textit{nach dort} und andere PPs wie [\textit{bis}[\textit{zum Strand}]]. 


Einige PPs\is{Präpositionalphrase (PP)} können vor der Präposition erweitert werden wie in \textit{kurz}/\textit{einen Meter vor der Brücke} oder \textit{gleich nach dem Essen}. Die vorangestellten Konstituenten sind formal nicht markiert. Sie hängen aber von der PP ab. Denn ohne die PP können sie nicht gebraucht werden (\ref{ex:*mod(PP)}), was umgekehrt nicht der Fall ist (\ref{ex:(mod)PP}). Es handelt sich um Konstituenten der PP.

\ea \label{ea:ModifPP}
\ea [] {Bleib [einen Meter vor der Brücke] stehen.}\label{ea:modPP}
\ex [*] {Bleib einen Meter stehen.}\label{ex:*mod(PP)}
\ex [] {Bleib [vor der Brücke] stehen.}\label{ex:(mod)PP}
\z 
\z 

Diese Erweiterung durch Maßangaben vor der Präposition ist nur bei lokalen PPs wie in (\ref{ea:modPP}) und bei temporalen PPs wie \textit{gleich nach dem Essen} möglich. 


\section{Subjunktionalphrase}\label{section:SubjunktP}

Analog zu den von einer Präposition regierten PP analysieren wir Nebensätze, die von einer Subjunktion eingeleitet werden, als Subjunktionalphrasen.\footnote{Für Subjunktionalphrasen verwenden wir keine Abkürzung, weil diese nicht gängig ist. Auch die Bezeichnungen sind unterschiedlich. \citet[241]{Pafel2011} bezeichnet sie als CP in Anlehnung an den englischen Begriff \textit{complementizer} für Subjunktion. \citet[367]{Schaefer2018_3Aufl} bezeichnet sie eingedeutscht als \textit{Komplementiererphrase} (KP).} Typische Subjunktionen sind \textit{dass}, \textit{ob}, \textit{als}, \textit{wenn}, \textit{weil}, \textit{obwohl}. Um die Wortart Subjunktion geht es in Abschnitt~\ref{subsubsection:Subjunktion}. Hier interessiert uns die von der Subjunktion regierte Subjunktionalphrase\is{Subjunktionalphrase}. In dem vorhergehenden Abschnitt haben wir gesehen, dass Präpositionen rechts von sich eine NP in einem ganz bestimmten Kasus regieren. Analog dazu regieren Subjunktionen rechts von sich eine finite VP mit Endstellung des Finitums (\ref{ea:Subjunktionalphrase}). 

\ea \label{ea:Subjunktionalphrase} dass das Kind draußen spielt
\z

Die Struktur der Subjunktionalphrase in (\ref{ea:Subjunktionalphrase}) stellen wir in der Abbildung~\ref{fig:SPAufbau} dar. Die Abkürzung S steht für Subjunktion.

\begin{figure}
  \centering
  \begin{forest}
    [Subjunktionalphrase, calign=child, calign child=1
      [S, tier=preterminal
        [dass]
      ]
      [ VPfin, tier=preterminal
        [das Kind draußen spielt, narroof]
      ]
      ]
  \end{forest}
  \caption{Aufbau einer typischen Subjunktionalphrase}
  \label{fig:SPAufbau}
\end{figure}



Insbesondere Subjunktionalphrasen mit \textit{bevor}, \textit{nachdem} und \textit{weil} können wie in (\ref{ea:SPerweitert}) -- ähnlich wie manche PPs -- durch eine vorangestellte Konstituente erweitert werden.  

\ea \label{ea:SPerweitert}
\ea \label{ea:Mod_tempSP}Suspendiert wurde der Beamte offenbar erst, [einige Wochen nachdem der Schmu aufgeflogen war].\footnote{Stern, 18.10.2018.}
\ex \label{ex:Mod_tempSP2} [Wenige Tage bevor die Castor-Behälter auf der Schiene nach Gorleben rollen], formieren sich die Atomkraft-Gegner im Wendland.\footnote{Braunschweiger Zeitung, 03.11.2010.}  
\ex \label{ex:Mod_kausSP} Das ist sehr ärgerlich, [besonders weil wir wirklich gleichwertig waren und die Chance auf den Sieg selbst verspielt haben], kommentiert der Team-Sprecher die Niederlage.\footnote{Braunschweiger Zeitung, 28.11.2005.}  
\z 
\z 

Diese Erweiterungen spezifizieren die temporale Bedeutung der ganzen Subjunktionalphrase: \textit{einige Wochen nachdem der Schmu aufgeflogen war} und \textit{wenige Tage bevor die Castor-Behälter auf der Schiene nach Gorleben rollen}. Es ist möglich, dass sie durch Konstruktionen wie \textit{einige Wochen, nach denen der Schmu aufgeflogen war} oder \textit{wenige Tage, vor denen die Castor-Behälter auf der Schiene nach Gorleben rollten} motiviert sind. Bei der kausalen Subjunktion \textit{weil} geht es eher um eine Art Graduierung wie bei \textit{besonders weil} oder eine Einschränkung des Wirklichkeitsbezugs wie \zB bei \textit{vielleicht weil} oder \textit{sicher weil}.

Abschließend sei auf die beiden Subjunktionen \textit{um} und \textit{ohne} hingewiesen. Diese regieren infinite VPs mit dem \textit{zu}-Infinitiv am Ende und bilden so wie in (\ref{ea:um_ohn_SP}) Subjunktionalphrasen.

\newpage
\ea \label{ea:um_ohn_SP}
\ea \label{ea:um_SP} um die Straße zu überqueren
\ex \label{ex:ohne_SP} ohne auf den Verkehr zu achten
\z 
\z 
 
Für (\ref{ea:um_SP}) stellen wir den Phrasenaufbau in der Abbildung~\ref{fig:SPAufbau_Inf} dar. 

 \begin{figure}
  \centering
  \begin{forest}
    [Subjunktionalphrase, calign=child, calign child=1
      [S, tier=preterminal
        [um]
      ]
      [ VPinf, tier=preterminal
        [die Straße zu überqueren, narroof]
      ]
      ]
  \end{forest}
  \caption{Aufbau einer Subjunktionalphrase mit infiniter \textit{zu}-VP}
  \label{fig:SPAufbau_Inf}
\end{figure}



\section{Adverbphrase}\label{section:Adverbphrase}

Typische\is{Adverbphrase (AdvP)} Adverbphrasen (AdvP) bestehen aus einem Adverb wie \textit{gestern} in (\ref{ea:AdvPtypisch}), das alleine eine Phrase bildet, was wir in Abbildung~\ref{fig:minimaleAdvP} darstellen.

\ea \label{ea:AdvPtypisch} Anna hat \textit{gestern} Paul getroffen.
\z 

\begin{figure}
\centering
\begin{forest}
    [AdvP
      [Adv, tier=preterminal
        [gestern]
      ]
    ]
\end{forest}
 \caption{minimale AdvP}
  \label{fig:minimaleAdvP}
\end{figure}

Anstelle des Adverbs kann auch eine Pro-Form den Kopf bilden. Statt \textit{gestern} könnte das Frage-Adverb \textit{wann} stehen. Adverbien haben gemeinsam, dass sie satzgliedfähig und nicht flektierbar sind. Um die Wortart Adverb geht es in Abschnitt~\ref{subsection:Adverb}.
 

Mitunter kommen zwei Adverbien vor wie \textit{hier oben} oder ein Adverb und eine PP wie \textit{oben auf dem Berg} in (\ref{ea:hieroben}). Jede der beiden Konstituenten kann wie in (\ref{ex:oben}) und (\ref{ex:aufdemBerg}) allein stehen. 

\ea \label{ea:AdvAdv/PP}
\ea \label{ea:hieroben} [Oben auf dem Berg] hat man eine gute Sicht.
\ex \label{ex:oben} [Oben] hat man eine gute Sicht.
\ex \label{ex:aufdemBerg} [Auf dem Berg] hat man eine gute Sicht.
\z 
\z 

Die formale Kennzeichnung und der Weglassbarkeitstest liefern keinen Aufschluss darüber, welche Konstituente von welcher abhängt, welche somit Teil der anderen wäre. Auch semantisch lässt sich nicht bestimmen, ob \textit{oben} eine nähere Bestimmung zu \textit{auf dem Berg} ist oder umgekehrt. Durch die gemeinsame Stellung im Vorfeld in (\ref{ea:hieroben}) sowie durch Erfragbarkeit (\textit{wo}) und Koordinierbarkeit (\textit{oben auf dem Berg und unten im Tal}) sollte angenommen werden, dass beide Konstituenten in lockerer Verbindung \textit{eine} größere Konstituente bilden. Ob diese koordinationsähnlich analysiert wird und/oder ad hoc als Einbettung der einen in die andere Konstituente, lassen wir offen.    


\section{Adjunktionalphrase}\label{section:Adjunktionalphrase}

Die in Vergleichsausdrücken gebrauchten Wörter \textit{als} und \textit{wie} bilden als Köpfe ebenfalls Phrasen. Dafür nehmen sie Nominalphrasen (\ref{ea:Adjunktor_NP}), Adjektivphrasen (\ref{ex:Adjunktor_AP}), Präpositionalphrasen (\ref{ex:Adjunktor_PP}) und Adverbphrasen (\ref{ex:Adjunktor_AdvP}) zu sich.

\ea\label{Adjunktorphrasen}
\ea\label{ea:Adjunktor_NP} Sie schläft [wie ein Murmeltier].
\ex\label{ex:Adjunktor_AP} Bei den meisten Lieferdiensten gelten die Arbeitsbedingungen [als sehr schlecht].
\ex\label{ex:Adjunktor_PP} [Wie in Deutschland] gibt es auch in Österreich eine allgemeine Schulpflicht von neun Jahren.
\ex\label{ex:Adjunktor_AdvP} Heute ist es wärmer  [als gestern].
\z 
\z      

Im Unterschied zu Präpositionen regieren \textit{als} und \textit{wie} in (\ref{Adjunktorphrasen}) keinen Kasus. Zwar können \textit{als} und \textit{wie} als Subjunktionen einen Nebensatz regieren und somit Subjunktionalphrasen (vgl. Abschnitt~\ref{section:SubjunktP}) bilden. Das ist jedoch in (\ref{Adjunktorphrasen}) nicht der Fall. In Abschnitt~\ref{subsubsection:Adjunktion} werden wir die Wortart von \textit{als} und \textit{wie} in (\ref{Adjunktorphrasen}) als \textit{Adjunktion} bestimmen. Folglich bezeichnen wir Phrasen, deren Köpfe Adjunktionen sind, als Adjunktionalphrasen.

Die Struktur der Adjunktionalphrase \textit{wie ein Murmeltier} in (\ref{ea:Adjunktor_NP}) stellen wir in der Abbildung~\ref{fig:Adjunktionalphrase Aufbau} dar.

\begin{figure}
  \centering
  \begin{forest}
    [Adjunktionalphrase, calign=child, calign child=1
      [Adjunktion, tier=preterminal
        [wie]
      ]
      [ NP, tier=preterminal
        [ein Murmeltier, narroof]
      ]
      ]
  \end{forest}
  \caption{Aufbau einer Adjunktionalphrase mit NP}
  \label{fig:Adjunktionalphrase Aufbau}
\end{figure}

Diese Phrasenstruktur gilt analog für Vergleichsausdrücke mit \textit{als} wie zum Beispiel (\textit{tiefer}) \textit{als ein Murmeltier}.



\section{Zusammenfassung Phrasen}\label{section:Zfsg_Phrasen}

Bevor wir uns im folgenden Kapitel~\ref{chap:Satzgliedanalyse} mit den möglichen syntaktischen Funktionen der hier vorgestellten Phrasen beschäftigen, fassen wir das Wichtigste zum Thema Phrasen in Stichpunkten zusammen. 

\begin{itemize}

\item Phrase: Zwischen der Wort- und der Satzebene gibt es die Ebene der Phrase. Häufig besteht eine Phrase aus mehreren Wörtern (Wortgruppe). Manche Wörter und Wortformen können allein eine Phrase bilden. 

\item Kopf: Jede Phrase besteht aus einem Kopf, der durch andere Einheiten ausgebaut werden kann oder muss. Die Wortart des Kopfes legt die Phrasenkategorie fest: finite Verbalphrase, infinite Verbalphrase, Nominalphrase, Adjektivphrase, Präpositionalphrase, Subjunktionalphrase, Adverbphrase, Adjunktionalphrase. 

\item Grammatische Beziehungen: Vom Kopf der Phrase können andere Konstituenten abhängen. Zwischen Kopf und den abhängigen Konstituenten bestehen drei weitere wesentliche Beziehungen: Rektion, Kongruenz und der Positionsbezug.

\item Rektion: Der Kopf legt aufgrund seiner lexikalischen und kategorialen Eigenschaften Formmerkmale der abhängigen Konstituenten fest.

\item Kongruenz: Bestimmte variable Merkmale der abhängigen Konstituenten stimmen mit den entsprechenden variablen Merkmalen des Kopfes überein.

\item Positionsbezug: Der Kopf hat an die von ihm abhängigen Konstituenten bestimmte Stellungsanforderungen.  

\end{itemize}

Der folgende letzte Abschnitt dieses Kapitels bietet Ihnen die Gelegenheit, durch Übungen zu lernen.

\section{Aufgaben zu Phrasen}\label{section:Aufgaben_Phrasen}

In den folgenden vier Aufgaben können Sie Ihr Wissen zum Thema Phrasen anwenden. Wir empfehlen, dass Sie die Aufgaben zuerst selbstständig lösen. Im Anschluss können Sie Ihre Lösungen mit unseren Lösungen und Kommentaren vergleichen.


\subsection{Phrasentypen}\label{subsection:A_Phrasentypen}

In diesem Kapitel (Kapitel \ref{chap:Phrasen}) haben Sie verschiedene Phrasentypen kennengelernt: finite und infinite Verbalphrasen (VP, vgl. Abschnitt~\ref{section:VP}), Nominalphrasen (NP, vgl. Abschnitt~\ref{section:NP}), Adjektivphrasen (AP, vgl. Abschnitt~\ref{section:AP}), Präpositionalphrasen (PP, vgl. Abschnitt~\ref{section:PP}), Subjunktionalphrasen (Abschnitt~\ref{section:SubjunktP}), Adverbphrasen (AdvP, vgl. Abschnitt~\ref{section:Adverbphrase}) und Adjunktionalphrasen (vgl. Abschnitt~\ref{section:Adjunktionalphrase}). Die in (\ref{ea:A_Phrasentypen}) folgenden, teils übersetzten Titel von Büchern, Geschichten und Filmen sind Phrasen.   

\ea \label{ea:A_Phrasentypen}
\ea \label{ea:A_AP} Dem Himmel so fern\footnote{Kajsa Ingemarsson. S. Fischer Verlage, 2014.} 
\ex \label{ex:A_PP} Auf der Suche nach der verlorenen Zeit\footnote{Marcel Proust. Knesebeck Verlag, 2010--2022.} 
\ex \label{ex:A_VPfin} Jeder stirbt für sich allein\footnote{Fallada, Hans. Aufbau-Verlag, [1947] 2011.}
\ex \label{ex:A_AdvP} Hier draußen\footnote{Martina Behm. dtv, 2025.} 
\ex \label{ex:A_NP1} Eines langen Tages Reise durch die Nacht\footnote{Eugene O'Neill. S. Fischer Verlage, 2002.}
\ex \label{ex:A_VPinf} Ins Holz gehen\footnote{Carlo Cassola. Kampa Verlag, 2024.}
\ex \label{ex:A_SP} Als die Zeit vom Himmel fiel\footnote{Mella Dumont. CreateSpace Independent Publishing Platform, 2015.}
\ex \label{ex:A_NP2} Hundert Jahre Einsamkeit\footnote{Gabriel García Márquez. Aufbau-Verlag, 1980.}
\ex \label{ex:A_AdjunktionalP} Wie ein einziger Tag\footnote{Nicholas Sparks. Heyne, 1996.}
\z 
\z 

Jeder Titel ist eine Phrase, in die andere Phrasen eingebettet sein können. Nennen Sie den Phrasentyp jedes Titels und das Wort, das der Kopf der gesamten Phrase ist. Begründen Sie den Kopfstatus.



Kommen wir nun zu den Lösungen: Die Wortgruppe \textit{dem Himmel so fern} in (\ref{ea:A_AP}) ist eine AP. Ihr Kopf ist das Adjektiv \textit{fern}. Der Kopf regiert den Dativ der NP \textit{dem Himmel}. 

Die Wortgruppe \textit{auf der Suche nach der verlorenen Zeit} in (\ref{ex:A_PP}) ist eine PP. Ihr Kopf ist die Präposition \textit{auf}, die nach einer NP im Dativ, \textit{der Suche nach der verlorenen Zeit}, verlangt. 

Die Wortgruppe \textit{jeder stirbt für sich allein} in (\ref{ex:A_VPfin}) ist ein Satz. Das Zentrum dieses Satzes ist das finite Verb \textit{stirbt}. Es verlangt nach einer NP im Nominativ, hier \textit{jeder}. Die PP \textit{für sich allein} ist nicht vom Verb gefordert, hängt aber von ihm ab. Sie ist, anders als das Verb, weglassbar. Die gesamte Wortgruppe ist eine finite VP. 

Die Wortgruppe \textit{hier draußen} in (\ref{ex:A_AdvP}) ist eine AdvP. Ob das Adverb \textit{hier} oder das Adverb \textit{draußen} der Kopf ist, ist unklar. In entsprechenden Kontexten könnte jedes der beiden Adverbien auch allein stehen, \zB \textit{hier spielen} oder \textit{draußen spielen}. Wir können sowohl \textit{hier} als nähere Bestimmung zu \textit{draußen} verstehen als auch umgekehrt, \textit{draußen} als nähere Bestimmung zu \textit{hier}.

Die archaisch anmutende Wortgruppe \textit{eines langen Tages Reise durch die Nacht} in (\ref{ex:A_NP1}) ist eine NP. Ihr Kopf ist das Nomen \textit{Reise}. Die vorangestellte genitivische NP \textit{eines langen Tages} und die nachgestellte PP \textit{durch die Nacht} hängen vom Kopfnomen ab und werden von ihm regiert.

Die Wortgruppe \textit{ins Holz gehen} in (\ref{ex:A_VPinf}) ist eine infinite VP. Ihr Kopf ist der Infinitiv \textit{gehen}. Er regiert die Richtungsergänzung \textit{ins Holz}.

Die Wortgruppe in (\ref{ex:A_SP}) \textit{als die Zeit vom Himmel fiel} ist ein Nebensatz, der von der Subjunktion \textit{als} eingeleitet ist. Diese Subjunktion ist der Kopf und verlangt nach einer finiten Verbalphrase mit Endstellung des Finitums. Der gesamte Nebensatz ist eine Subjunktionalphrase.

Die Wortgruppe \textit{hundert Jahre Einsamkeit} in (\ref{ex:A_NP2}) ist eine NP. Ihr Kopf ist das Nomen \textit{Jahre}. Sowohl \textit{hundert} als auch \textit{Einsamkeit} hängen vom Kopf ab. Die ganze NP steht im Plural, weil der Kopf \textit{Jahre} im Plural steht.

Die Wortgruppe \textit{wie ein einziger Tag} in (\ref{ex:A_AdjunktionalP}) ist eine Adjunktionalphrase. Ihr Kopf ist die Adjunktion \textit{wie}. Der Kopf \textit{wie} legt als Positionsbezug fest, dass die NP \textit{ein einziger Tag} rechts von ihm steht.




\subsection{Beziehungen in Phrasen}\label{subsection:A_Beziehungen_Phrasen}

Innerhalb von Phrasen gibt es drei wesentliche Beziehungen zwischen dem Kopf und den zu ihm gehörenden Wörtern und Phrasen, nämlich die Rektion, die Kongruenz und den Positionsbezug. Diese Beziehungen bestehen auch in der in (\ref{ea:A_Beziehungen_Phrasen}) folgenden komplexen Phrase.   

\ea \label{ea:A_Beziehungen_Phrasen} der Teufel mit den drei goldenen Haaren\footnote{Titel eines Märchens aus: Jacob und Wilhelm Grimm: \textit{Kinder- und Hausmärchen}. [KHM 29]. Erstdruck 1812. Berlin: Europäischer Literaturverlag, 2015.}
\z 

Nennen Sie jeweils zwei Beispiele für Rektion, für Kongruenz und für den Positionsbezug innerhalb dieser komplexen Phrase.



Kommen wir nun zu den Lösungen: In (\ref{ea:A_Beziehungen_Phrasen}) regiert das Kopfnomen \textit{Teufel} das maskuline Genus des Artikels \textit{der}. Die Präposition \textit{mit} regiert den Dativ der NP \textit{den drei goldenen Haaren}. Innerhalb der NP kongruieren die dativischen Pluralformen \textit{den}, \textit{goldenen} und \textit{Haaren}. Das Kopfnomen \textit{Haaren} regiert das Genus Neutrum, was jedoch im Plural nicht sichtbar wird. Der Positionsbezug zeigt sich in der Stellung der PP \textit{mit den drei goldenen Haaren} nach dem Kopfnomen \textit{Teufel}, im sogenannten nominalen Nachfeld. Der Artikel \textit{der} bildet die linke Nominalklammer und steht, wenn das Mittelfeld unbesetzt ist, direkt vor dem Nomen \textit{Teufel}, das die rechte Nominalklammer bildet. Um Positionsbezug handelt es sich auch innerhalb der PP bei der Spitzenstellung der Präposition \textit{mit} und der Nachstellung der von ihr regierten NP \textit{den drei goldenen Haaren}. Auch die Stellung innerhalb der NP ist geregelt. Der Artikel \textit{den} bildet die linke Nominalklammer und das Kopfnomen \textit{Haaren} die rechte. Zwischen beiden, im sogenannten Mittelfeld, stehen \textit{drei goldenen}. Das Nachfeld ist unbesetzt.  
 

\subsection{Kopf und Nicht-Kopf}\label{subsection:A_Nicht_Kopf_Phrasen}

Jede Phrase hat einen Kopf, der die Merkmale der gesamten Phrase festlegt. In den folgenden Beispielen ist jeweils ein Wort kursiv gesetzt, das als Kopf einer Phrase (die selbst wieder Bestandteil einer anderen Phrase ist) fungiert.

\ea \label{ea:A_Phrasen_Kopf}
\ea \label{ea:A_Kopf1} Eine Geschichte \textit{von} Liebe und Finsternis\footnote{Amos Oz. Suhrkamp Verlag, 2004.}  
\ex \label{ex:A_Kopf2} Ziemlich \textit{beste} Freunde\footnote{Ins Deutsche übersetzter Filmtitel von Olivier Nakache und Éric Toledano.}
\ex \label{ex:A_Kopf3} Ein Tag im \textit{Leben} des Iwan Denissowitsch\footnote{Alexander Solschenizyn. Knaur, 1963.}
\ex \label{ex:A_Kopf4} Der Tag, \textit{als} meine Frau einen Mann fand\footnote{Sibylle Berg. Hanser, 2015.}
\z 
\z 

Nennen Sie die Phrasen, die von dem kursiv gesetzten Kopf gebildet werden, und ordnen Sie sie einem Phrasentyp zu.



Kommen wir nun zu den Lösungen: In (\ref{ea:A_Kopf1}) ist die Präposition \textit{von} der Kopf der PP \textit{von Liebe und Finsternis}. Das Adjektiv \textit{beste} in (\ref{ex:A_Kopf2}) ist der Kopf der AP \textit{ziemlich beste}. Das Nomen \textit{Leben} in (\ref{ex:A_Kopf3}) ist der Kopf der NP \textit{Leben des Iwan Denissowitsch}. In (\ref{ex:A_Kopf4}) ist die Subjunktion \textit{als} der Kopf der Subjunktionalphrase \textit{als meine Frau einen Mann fand}.


\subsection{Felder in der Nominalphrase}\label{subsection:A_Felder_NP}

Analog zur Satzklammer haben wir von einer Nominalklammer (vgl. Abschnitt~\ref{section:NP}) gesprochen. Diese wird durch den Artikel als linke Nominalklammer und das Kopfnomen als rechte Nominalklammer gebildet. In dem Mittelfeld zwischen beiden können APs stehen, deren Kopf ein flektiertes Adjektiv ist. Genitivische NPs, PPs und Relativsätze können nach dem Kopfnomen im sogenannten nominalen Nachfeld stehen. Analysieren Sie die komplexen Nominalphrasen in (\ref{ea:A_komplexeNP}), die von dem jeweils kursiv gesetzten Kopf gebildet werden. 

\ea \label{ea:A_komplexeNP}
\ea \label{ea:A_komplexeNP1} Immerhin ist der Stiftungsakt ein sichtbares Zeichen, dass der jahrelange \textit{Streit} um den Wiederaufbau der Garnisonkirche beendet ist.\footnote{Hannoversche Allgemeine, 21.06.2008.}
\ex \label{ex:A_komplexeNP2} Das Angeln gefällt auch Wolodja, und so gewinnt der Agent schnell das Vertrauen des schwatzhaften und auf seinen berühmten Vater stolzen \textit{Knaben}.\footnote{Die Zeit, 11.03.1954.}
\ex \label{ex:A_komplexeNP3} Teuer dürfte für einen 44-Jährigen der \textit{Racheakt} im Streit um einen entgangenen Parkplatz am Berliner Platz werden.\footnote{Braunschweiger Zeitung, 09.02.2010.}
\z 
\z 

Nennen Sie zunächst die gesamte Nominalphrase. Geben Sie dann die Besetzung der Felder an. Bestimmen Sie restlos alle Phrasentypen und ihre Bestandteile, die in der Nominalphrase enthalten sind.



Kommen wir nun zu den Lösungen: In (\ref{ea:A_komplexeNP1}) ist \textit{Streit} der Kopf der NP \textit{der jahrelange Streit um den Wiederaufbau der Garnisonkirche}. Der Artikel \textit{der} bildet die linke Nominalklammer und das Kopfnomen \textit{Streit} die rechte. Im Mittelfeld steht \textit{jahrelange}. Im Nachfeld steht \textit{um den Wiederaufbau der Garnisonkirche}. Das Adjektiv \textit{jahrelange} bildet eine AP, die Bestandteil der NP ist. Bei \textit{um den Wiederaufbau der Garnisionskirche} handelt es sich um eine PP, die aus der Präposition \textit{um} und der komplexen NP \textit{den Wiederaufbau der Garnisionskirche} besteht. Die NP besteht aus dem Artikel \textit{den} und dem Kopfnomen \textit{Wiederaufbau}, dem die genitivische NP \textit{der Garnisionskirche} untergeordnet ist.

In (\ref{ex:A_komplexeNP2}) ist \textit{Knaben} der Kopf der NP \textit{des schwatzhaften und auf seinen berühmten Vater stolzen Knaben}. Der Artikel \textit{des} bildet die linke Nominalklammer und das Kopfnomen \textit{Knaben} die rechte. Im Mittelfeld steht \textit{schwatzhaften und auf seinen berühmten Vater stolzen}. Das Nachfeld ist unbesetzt. Das Adjektiv \textit{schwatzhaften} bildet eine AP, die durch die Konjunktion \textit{und} mit der AP \textit{auf seinen berühmten Vater stolze} zu einer AP verbunden ist. Der Kopf der AP \textit{auf seinen berühmten Vater stolze} ist das Adjektiv \textit{stolze}, das die PP \textit{auf seinen berühmten Vater} regiert. Diese besteht aus der Präposition \textit{auf} und der NP \textit{seinen berühmten Vater}. Der Kopf der NP ist \textit{Vater}. Zwischen dem Artikel \textit{seinen} und dem Kopfnomen steht das Adjektiv \textit{berühmten}, das eine AP bildet.  

In (\ref{ex:A_komplexeNP3}) ist \textit{Racheakt} der Kopf der NP \textit{der Racheakt im Streit um einen entgangenen Parkplatz am Berliner Platz}. Der Artikel \textit{der} bildet die linke Nominalklammer und das Kopfnomen \textit{Racheakt} die rechte. Das Mittelfeld ist unbesetzt. Im Nachfeld steht \textit{im Streit um einen entgangenen Parkplatz am Berliner Platz}. Diese Wortgruppe ist eine PP mit der Präposition \textit{im} als Kopf und der komplexen NP \textit{Streit um einen entgangenen Parkplatz am Berliner Platz}. Ihr Kopf ist das Nomen \textit{Streit}, das die PP \textit{um einen entgangenen Parkplatz am Berliner Platz} regiert. Ihr Kopf ist die Präposition \textit{um}, die die komplexe NP \textit{einen entgangenen Parkplatz am Berliner Platz} regiert. Der Kopf dieser NP ist \textit{Parkplatz}. Zwischen dem Artikel \textit{einen} und dem Kopfnomen \textit{Parkplatz} steht das Adjektiv \textit{entgangenen}, das eine AP bildet. Im Nachfeld des Kopfnomens \textit{Parkplatz} steht die PP \textit{am Berliner Platz}. Ihr Kopf ist die Präposition \textit{am}, die die NP \textit{Berliner Platz} regiert. Ihr Kopf ist das Nomen \textit{Platz}. Das Adjektiv \textit{Berliner} bildet eine AP, die dem Kopfnomen untergeordnet ist. Alternativ können wir die PP \textit{am Berliner Platz} auch als Ort des Streites verstehen. Dann ist sie nicht \textit{Parkplatz} als Kopf untergeordnet, sondern \textit{Streit}.
